\documentclass[10pt,letterpaper]{article}

\usepackage[T1]{fontenc}
\usepackage[utf8]{inputenc}
\usepackage{lmodern}
\usepackage[margin=0.68in]{geometry}
\usepackage{amsmath,amssymb,mathrsfs}
\usepackage{xcolor}
\usepackage{microtype}
\usepackage[hidelinks]{hyperref}
\hypersetup{
  pdftitle={The claimed complex structure on S6: manuscript and proof-reconstruction overview},
  pdfauthor={Codex mathematical reconstruction with ChatGPT Pro review records},
  pdfsubject={Equation-first mathematical account of the candidate, topology, Hodge and Bott-Chern calculations, analytic Picard completion, toric Kummer covers, cusp limits, cone-target sphere maps, split-Cartan SU(r) gauge systems, rank-two and rank-three SU(3) Gram quotients, the Lorentz observation morphism, Hermitian determinant normal slice, Weyl-spinor-scalar splitting, colour-invariant momentum, Klein-Gordon and compact gauge-scalar models, the S8 obstruction, divisor-defect and normal-holonomy lattices, the SU(2)-by-SU(3) tensor subgroup of SU(6), the Erdos-Straus-Koide incidence bridge, and the CDP20 consequence},
  pdfkeywords={S6, S8, complex threefold, nearby cycles, Leray spectral sequence, CDP20, toric Kummer cover, cusp degeneration, pointed limit, smash product, SU(r) flux, SU(3) Gram quotient, SU(6), normal holonomy, Hermitian Jordan cone, Hermitian normal slice, Lorentz observation morphism, spin representation, Klein-Gordon, gauge-scalar model, coercive spectral transfer, colour-invariant momentum, determinant phase, defect lattice, Erdos-Straus, Koide cone, proof audit}
}
\setlength{\parindent}{1.1em}
\setlength{\parskip}{2.6pt}
\setlength{\abovedisplayskip}{4.2pt plus 1.2pt minus 1.5pt}
\setlength{\belowdisplayskip}{4.2pt plus 1.2pt minus 1.5pt}
\setlength{\emergencystretch}{2em}
\pagestyle{plain}
\frenchspacing
\newcommand{\overviewsection}[1]{%
  \noindent{\large\bfseries #1}\par\vspace{1.5pt}}
\DeclareMathSizes{9.4}{9.4}{7}{5}

\begin{document}
\ifdefined\SIXOverviewEmbedded
\else
\begin{center}
{\Large\bfseries The Claimed Complex Structure on $S^6$}\par
\vspace{1.5pt}
{\large What the manuscript and the proof reconstruction establish}\par
\vspace{2pt}
{\small A standalone mathematical orientation to the canonical workbench}\par
\end{center}
\vspace{-4pt}
\fi
\fontsize{9.4pt}{10.65pt}\selectfont

\noindent\textbf{Document map.}
This file is the equation-first orientation to the cumulative
\path{workbench.tex}.  Readers following one theorem lane can instead use the
standalone drafts below; the cumulative workbench remains the source of record
for their shared hypotheses and comparison maps.
\begin{flushleft}\small
\path{koide_arithmetic_preprint.tex}: exact Koide arithmetic.\\
\path{s6_spectral_degeneration_preprint.tex}: cusp spectral limits.\\
\path{s6_compact_link_gaussian_preprint.tex}: compact-link Gaussian model.\\
\path{s6_cone_sphere_routes_preprint.tex}: sphere-carrier topology.\\
\path{s6_erdos_straus_koide_bridge_preprint.tex}: marked Erd\H{o}s--Straus--$S^6$
incidence, the modulo-$840$ lattice and exact shell maps.\\
\path{s6_koide_es_divisor_extension_audit.tex}: divisor-channel audit and two distinct arithmetic maps.\\
\path{s6_higher_rung_su3_gram_cone_preprint.tex}: rank-three Gram and toroidal-successor calculations.\\
\path{s6_worldline_presymplectic_reduction_preprint.tex}: exact massive,
massless, and hard-threshold worldline reduction.
\\\path{s6_quaternionic_4441_reconciliation.tex}: common quaternionic slice,
positive fibres, retained Jordan product, omitted-coordinate dynamics,
exact rotating coupling fibres, the smooth literal scalar--Gauss domain,
and generic spatial exclusion.
\\\path{s6_homogeneous_quantum_confinement.tex}: complete homogeneous
gauge--scalar quantization, Gauss law, and positive finite-model gap.
\\\path{s6_higher_rung_24d_preprint.tex}: higher torus wrapping,
octonionic determinant boundary, triality flag, exact mod--12/mod--24 maps,
and the shared residue-lattice/Weyl-quotient continuation.
\end{flushleft}

\noindent\textbf{6 September: explicit spectral and umbral maps, and the original family's parameter identifications.}
For the literal common lattice $\mathcal I=\{\sqrt2(a,a,a):a\in D_4\}$,
write $a=a_0+v$, $t^2=|v|^2$, $Q=a_0^2+t^2$ and
$P=a_0(a_0^2-3t^2)$. Its complete cubic image is
\[
 \tau(\mathcal I)=2\sqrt2\bigl(18\mathbb Z\cup(6\mathbb Z+2)
                                      \cup(6\mathbb Z-2)\bigr).
\]
Three explicit integral affine lines give a section of this value map.
The scalar-retaining Albert matrix $X=\mu I+Y(a)$, with cyclic
off-diagonal entries $\sqrt2a$, has invariants and labelled spectrum
\[
\begin{aligned}
 (\operatorname{tr}X,\operatorname{tr}X^2,\det X)
   &=(3\mu,3\mu^2+12Q,\mu^3-6\mu Q+4\sqrt2P),\\
 \operatorname{Spec}Y
   &=(2\sqrt2a_0,-\sqrt2a_0-\sqrt6t,-\sqrt2a_0+\sqrt6t).
\end{aligned}
\]
For $t>0$, the explicit labelled eigenprojector satisfies
$(\Pi_1)_{12}=-1/6-\sqrt3(v/t)/6$; thus the frame and spectrum
recover $a$, not only its norm and cubic. The full discriminant is
$864t^2(3a_0^2-t^2)^2$. The full Cartan family retains the scalar
$\mu$, radius $r$, and angle $\theta$, with $r^2=12Q$ and
$r^3\cos3\theta=24\sqrt3P$ on this subfamily.

For the same marked root system $A_5^4D_4$, the
Cheng--Duncan--Harvey prescription gives
$S^X=(5I_{12}+J)S_6$, where $J_{rs}=\mathbf1_{s=5r\ (12)}$ and
$S_{6,r}=\sum_{k\equiv r\ (12)}kq^{k^2/24}$.
The coefficient map to the retained Fricke function is
\[
 T_6=\frac{\eta(\tau)^5\eta(3\tau)}{\eta(2\tau)\eta(6\tau)^5},
 \qquad f^X=-q\frac{d}{dq}\log T_6,
 \qquad T_6\!\left(-\frac1{6\tau}\right)=\frac{72}{T_6(\tau)}.
\]
All twelve theta indices, the multiplier determinant, the original
glue-group action and its twisted shadows are computed in the proofs.
This is an explicit use and comparison of established umbral machinery,
not a claim to a new moonshine module.

In the original analytic family $\beta_c=\beta^{\rm part}+c$,
$\operatorname{Im}c<-M$, the actual global identifications are exactly
$X_c\simeq X_{c'}\Longleftrightarrow c'-c\in2\mathbb Z$.
The middle lattices agree for every integer shift; the multiplicity-four
filling retains an order-two obstruction to odd shifts. Even shifts have
explicit finite-filling and toric cusp maps and a proved global additive
gluing correction. The resulting parameter quotient is
$w=e^{-\pi ic}$, $0<|w|<e^{-\pi M}$.
This classifies the retained analytic family without using the claimed
identification of its total space with $S^6$.


\noindent\textbf{6 September: cubic arithmetic, integral lattice maps, and the canonical ring.}
For the retained $A_5^4\perp D_4$ Niemeier lattice $N$ and its
original index-six Leech neighbour $\Lambda$, fix the Cayley--Dickson
frame $(1,i,j,k,\ell,i\ell,j\ell,k\ell)$ used in the complete cubic
proof. Put $\tau(q)=\operatorname{Re}((q_1q_2)q_3)$ and
$F=\tau\circ\mathcal L$ on the original Niemeier span,
and $\mathcal V_F(L)=\operatorname{span}_{\mathbb Z}\{F(y):y\in L\}$.
The exact additive groups are
\[
\begin{aligned}
 \mathcal V_F(N)&=\tfrac1{27}\mathbb Z\oplus
       \tfrac{\sqrt2}{54}\mathbb Z\oplus\tfrac{\sqrt3}{18}\mathbb Z,\\
 \mathcal V_F(\Lambda)&=\tfrac1{27}\mathbb Z\oplus
       \tfrac{\sqrt2}{216}\mathbb Z\oplus\tfrac{\sqrt3}{54}\mathbb Z,
 \qquad \mathcal V_F(\Lambda)/\mathcal V_F(N)\cong\mathbb Z/12.
\end{aligned}
\]
The Albert determinant is $\lambda^3-\lambda(y,y)+2F(y)$.
The quotient concerns the displayed value-generated groups in this fixed
frame, not an inclusion $N\subset\Lambda$. The actual residue
$r=\rho(361)$ cycles three Leech neighbours $\Lambda_j=r^j\Lambda$.
Their pairwise intersection is one integral lattice $J$, with
$\det J=81$ and $[\Lambda_j:J]=9$. The fixed fourth isotropic plane
in $J^\vee/J$ gives $M$ with exactly $72$ roots of type $A_2^{12}$,
$M\cap\Lambda_j=J$, and $[M:J]=9$.
\[
 \mathcal V_F(J)=\mathcal V_F(M)
 =\tfrac1{27}\mathbb Z\oplus\tfrac{\sqrt2}{216}\mathbb Z
                         \oplus\tfrac{\sqrt3}{18}\mathbb Z.
\]
The successive value-group inclusions have matrices
$\operatorname{diag}(1,4,1)$ and $\operatorname{diag}(1,1,3)$,
giving quotients $\mathbb Z/4$ and $\mathbb Z/3$.
Both containments and their reverse generator certificates are calculated
from the full cubic polarization, not inferred from lattice indices.

Wilson's construction \href{https://doi.org/10.1016/j.jalgebra.2009.03.021}{(R.~A.~Wilson, 2009)}, transported by the explicit
isometry $\iota(x)=x/\sqrt2$, gives $\Lambda_{\rm cyc}$ in the
original triality metric. Its cyclic fixed Gram matrix is $3G_{E_8}$;
its orthogonal complement is explicitly isometric to $A_2\otimes E_8$,
and their sum has index $3^8$, with quotient $\mathbb F_3^8$.
The directions remain $c(q_1,q_2,q_3)=(q_3,q_1,q_2)$ and
$c_{\rm W}=c^{-1}$. Its literal intersections are
\[
 \Lambda_{\rm cyc}\cap\mathcal L(N)
 =\Lambda_{\rm cyc}\cap\mathcal L(\Lambda)
 =\Delta:=\{\sqrt2(a,a,a):a\in D_{\mathbb H}\},
\]
where $D_{\mathbb H}$ is the even-coordinate-sum integer lattice in
the first four stated Cayley coordinates. Thus $\Delta$ has rank $4$,
Gram matrix $6G_{D_4}$, determinant $5184$, minimum squared norm $12$,
and value-generated subgroup $4\sqrt2\mathbb Z$ for $\tau|_\Delta$, or
$8\sqrt2\mathbb Z$ for $2\tau|_\Delta$. This is not the actual image:
$12\sqrt2$ is not attained by $\tau$ on $\Delta$.
The preceding proofs retain the triality products and spectral boundary
graph, the flag's degree-$+1$ collapse, and the retained cusp class's
nontrivial three-plane reduction; these exact maps do not supply an
interacting quantum mass-gap result.

For the original analytic fibration $f:X\to\mathbb P^1$, the full
canonical-divisor proof gives
$\omega_X\simeq\mathcal O_X(-2S_2)$ and the graded anticanonical ring
$\bigoplus_{m\ge0}H^0(X,\omega_X^{-m})\cong\mathbb C[U,V]$,
with $\deg U=1$, $\deg V=2$. The complete system $|-2K_X|$
recovers $f$ with the specified projective coordinates. This calculation
uses the original analytic charts and finite fillings, not an
identification of $X$ with $S^6$, and draws no sphere-topology conclusion.

The retained constant $\beta_c=\beta^{\rm part}+c$ now gives a proper
holomorphic deformation over $\operatorname{Im}c<-M$, locally on
relatively compact parameter discs. On a fixed regular fibre set
$T=\operatorname{Im}\tau$, $m=\operatorname{Im}\mu$,
$D=\operatorname{Im}\beta_{c_*}-6m^2/T<0$ and $\delta=c-c_*$.
In the plus-graph convention its exact Beltrami matrix is
\[
 \nu_\delta=\frac{\delta}{2iD+\delta}
                 \begin{pmatrix}0&0\\-m/T&1\end{pmatrix}.
\]
The connecting Kodaira--Spencer class is
$-[\partial_{\zeta_2}\otimes
(d\bar\zeta_2-(m/T)d\bar\zeta_1)/(2iD)]\ne0$.
The intrinsic pencil recovers $f$ also over the dual numbers; its three
fixed critical values prohibit an infinitesimal projective base change.
Consequently this deformation of the whole $X_c$ is nontrivial,
not only its marked regular-fibre deformation. The proof uses no sphere
recognition or CDP conclusion.


\noindent\textbf{5 September: two corrections to the S6 proof text.}
For a compact Cartier surface $W$ in a smooth complex threefold, the
duality map is constructed by the analytic trace pairing with
$K_W^\bullet=\omega_W[2]$; matching outer terms of exact sequences
does not construct its middle arrow. For the logarithmic gluing
$\zeta\mapsto\zeta+\alpha(s)\Pi(s)v$, the exact real translation is
\[
 a(s)=\operatorname{Re}\alpha(s)\,v+
       \operatorname{Im}\alpha(s)\,J_sv,\qquad
 J_s=\Pi(s)^{-1}i\Pi(s),\quad \alpha(s)=\log s/(2\pi i).
\]
On a full clockwise covering circle, $a(1)-a(0)=-v$, so the mixed
integral is $-E(v,\lambda)$ for each marked fibre circle $\lambda$.
The period covariance gives $J_{gs}A=AJ_s$ and
$a(gs)=Aa(s)-v/m$. These replace the incorrect real-scalar expression
while preserving the finite-norm identity and based meridian relation.
The workbench supplies the complete corrected proofs and the exact
full-fibre map
$(R^2f_*E)_p/\mathfrak m_p(R^2f_*E)_p\cong H^2(W,E|_W)$.
This recheck concerns those maps, not a fresh certification of every
global input or a theorem-level counterexample to CDP20.

\noindent\textbf{5 September: the exact period half-plane and cusp identifications.}
For the cusp-regular particular solution set
$D_{\rm part}=\operatorname{Im}\beta^{\rm part}
-6(\operatorname{Im}\mu)^2/\operatorname{Im}\tau$ and
$M=\max_{B\setminus\{p_0\}}D_{\rm part}$. This finite maximum is
attained. The exact constants for which $\beta=\beta^{\rm part}+c_0$
gives a real period isomorphism everywhere are
\[
 \{c_0\in\mathbb C:\operatorname{Im}c_0<-M\}.
\]
Without the cusp extension, $\widetilde\beta=\beta+i\,t\circ\pi$
obeys the same affine laws but has
$\widetilde D=D+\operatorname{Re}t\to+\infty$ along the positive
real base ray. The proof is the proposition ``exact admissible period
constants'' and its following remark in \path{candidate_geometry.tex}.
The cusp proof in that same source now states the restricted charts
$\mathcal U_\sigma(S,\epsilon)=\mathcal U_\sigma(S)\cap Y_\epsilon$,
uses a bound for the full $C(t_c)$ on a fixed closed subdisc, and retains
the closed-polydisc condition $\theta_i^\sigma(y)\ge0$ for torus points.
The explicit opposite-edge maps below retain all entries of $C(0)$.
These period and local-quotient refinements, together with the bounded
finite-filling recheck, do not add a sphere-identification or CDP theorem.

\noindent\textbf{5 September: the 24-dimensional continuation.}
The rank-three octonionic Jordan algebra provides two explicit geometric
objects. For $\lambda\in\mathbb R$, $q_1,q_2,q_3\in\mathbb O$, set
\[
 Q(\lambda;q)=
 \begin{pmatrix}
 \lambda&q_3&\bar q_2\\
 \bar q_3&\lambda&q_1\\
 q_2&\bar q_1&\lambda
 \end{pmatrix},\qquad
 \det Q=\lambda^3-\lambda\sum_{i=1}^3|q_i|^2
             +2\operatorname{Re}((q_1q_2)q_3).
\]
Its positive cone $\mathcal C=\{Q\succeq0\}$ lies in the displayed
25-dimensional real slice. The slice $\lambda=1$ is a compact convex
24-ball; radial coordinates give an explicit homeomorphism of the cone
and its determinant-zero boundary to a closed half-space and its boundary.
One-point compactification therefore gives
\[
 (\mathcal C^+,(\partial\mathcal C)^+)\cong(B^{25},S^{24}).
\]
The second object, the octonionic flag $\operatorname{Fl}(\mathbb O)
=F_4/\operatorname{Spin}(8)$, has three rank-eight tangent summands
$8_v,8_s^+,8_s^-$ and integral cohomology
\[
 \frac{\mathbb Z[\beta_1,\beta_2]}
 {(\beta_1^2-\beta_1\beta_2+\beta_2^2,
       \beta_1\beta_2(\beta_1-\beta_2))},\qquad |\beta_i|=8.
\]
Its six cells have dimensions $0,8,8,16,16,24$ and its Euler class is
$6\beta_1^2\beta_2$. These known flag results are attributed to
\href{https://arxiv.org/abs/0809.4318}{Mare--Willems}; the draft writes
the comparison with the retained toric and determinant constructions.
The six flag fixed points and the six toric zero-strata admit an explicit
$S_3$-equivariant bijection. That bijection alone supplies no map on their
positive-dimensional strata. Both the determinant-boundary sphere and
the fourth smash of a chosen based $S^6$ carrier are $S^{24}$, with a
displayed choice-dependent homeomorphism. The flag has degree-one
collapse maps onto either sphere; its nonzero $H^8$ prevents homotopy
equivalence with a sphere.

The mod--12/mod--24 relation has explicit maps. In the retained peripheral
representation,
\[
 \chi:SL_2(\mathbb Z)\longrightarrow\mathbb Z/12,\qquad
 (\chi(C_0),\chi(C_1),\chi(C'_2))=(-1,4,-3).
\]
Every homomorphism $SL_2(\mathbb Z)\to\mathbb Z/24$ has even image.
The embedding $j(a)=2a$ exists, while no homomorphism lifts $\chi$
under ordinary reduction $\mathbb Z/24\to\mathbb Z/12$.
For a based cusp map $g:S^6\to S^3$, the suspension calculation gives
\[
 E^6:\pi_6(S^3)\cong\mathbb Z/12
       \hookrightarrow\pi_{12}(S^9)\cong\mathbb Z/24,
 \qquad E^6(a\alpha)=2a\nu_9.
\]
Thus $M\mapsto\chi(M)[g]$ has suspended value $2a\chi(M)\nu_9$ if
$[g]=a\alpha$, and image order $12/\gcd(a,12)$.
The actual coefficient is now computed: $a=0$, so this composite is zero.
For the original periods $Z(s)=sB_0+C(t)$, choose
$u^0=(-1/4,1/4)$ and the original chart $t=z_0z_1z_2$.
With $h_j(t)=\exp(-2\pi i(C(t)u^0)_j)$ the exact invertible map
\[
 (z_0,z_1,z_2)\longmapsto
 (z_0/(h_1h_2),h_1z_1,h_2z_2)=(w_0,w_1,w_2)
\]
preserves $t$ and puts the complete regular inverse image at
$(|w_0|,|w_1|,|w_2|)=(r^{1/2},r^{1/4},r^{1/4})$.
Writing $A=\log(r)B_0-2\pi\operatorname{Im}C(t)$, its actual
normal frame is $NJH^{-1}$, where $N$ consists of the three radial
normals, $H$ is the target-coordinate derivative, and
\[
 J=\begin{pmatrix}
 1/2&-A_{11}-A_{21}&-A_{12}-A_{22}\\
 1/4&A_{11}&A_{12}\\1/4&A_{21}&A_{22}
 \end{pmatrix},\qquad \det J=\det A.
\]
The bound $\|Av\|\geq(|\log r|-\|2\pi\operatorname{Im}C(t)\|)\|v\|$
keeps the framing homotopy invertible. The resulting constant change of
radial frame extends over a collared $D^2\times T^2$ inside the same
chart times an interval. Pontryagin--Thom therefore gives $g\simeq *$
for the original smooth source, without using $X\cong S^6$.
The complete chart, inverse, derivative, cap, and based-homotopy proof
is in the higher-rung draft, subsection \emph{The original framed cusp
coefficient is zero}, and in the cone/sphere draft.
The full draft gives the maps, the suspension proof, the source citations,
and the surviving local higher-torus morphisms. A global higher torus
completion with the required exceptional structure is still open.

\noindent\textbf{Retaining the angle detects topology that the preceding
projection loses.}
For the same continuous cusp map
$q_2:X\to\mathcal A=\operatorname{Susp}(S^1\times\overline T^2)$,
collapse the coordinate two-skeleton of $S^1\times\overline T^2$
to obtain $c_3:T^3\to S^3$ and set
$H([r,z])=[r,c_3(z)]$ after $q_2$. This defines
$H=p_3q_2:X\to S^4$. Keep the same original periods and coordinate
map above, but now fix $t_0=r_0e^{i\pi}$ as well as
$u^0=(-1/4,1/4)$. The complete regular inverse image is
\[
 w(a,b)=\bigl(r_0^{1/2}e^{i(\pi-a-b)},
              r_0^{1/4}e^{ia},r_0^{1/4}e^{ib}\bigr),
 \qquad (a,b)\in(\mathbb R/2\pi\mathbb Z)^2.
\]
The full derivative retains the period correction: with
$v=(\log|w_1|,\log|w_2|)$, $\ell=\log|t|$, and
$p_0=(1/4,1/4)$, the identity
$u=u^0+A(t)^{-1}(v-\ell p_0)$ gives, on this torus,
$du=A(t_0)^{-1}(dv-p_0d\ell)$. The term differentiating
$A(t)^{-1}$ vanishes there because its multiplying vector is exactly
zero. The resulting complete normal-plus-tangent frame has comparison
matrix $D(a,b)G_0$, where
\[
 D(a,b)=\operatorname{diag}(e^{-i(a+b)},e^{ia},e^{ib})\in SU(3)
\]
and $G_0$ is one constant invertible real matrix, with its orientation
sign calculated in the proof. An explicit periodic contraction of $D$
identifies the induced stable tangent framing with the Lie framing of
$T^2$. Its spin quadratic form is $q(x,y)=x+xy+y$ over
$\mathbb F_2$, with Arf invariant one. Pontryagin--Thom and the
injective stable-range suspension therefore give
\[
 [H]=\eta_4\eta_5\ne0\in\pi_6(S^4)\cong\mathbb Z/2,
 \qquad q_2\not\simeq_*0.
\]
Here $X$ has the fixed sphere identification used in the reconstruction.
This is not a new proof of that identification. Postcomposition by
$p_3$ gives the exact quotient
$\pi_6(\mathcal A)/\ker(p_3)_*\cong\mathbb Z/2$ and sends
the original $q_2$ to its nonzero element. In particular every odd
multiple of $[q_2]$ is nonzero; its full order is not determined here.

This retained class also supplies a nonzero mod--24 map. Fix the
complex Hopf map $\eta_2:S^3\to S^2$. The ordered smash gives
\[
 \mu_\eta:\pi_6(S^4)\longrightarrow\pi_9(S^6),
 \qquad [x]\longmapsto[\eta_2\wedge x],\qquad
 \mu_\eta([H])=\eta^3=12\nu_6\ne0.
\]
The equality uses the classical integral stable relation
$\eta^3=12\nu$, not $4\nu$ from its two-primary quotient
(\href{https://pi.math.cornell.edu/~hatcher/AT/AT.pdf}{Hatcher,
printed p.~385}). Three suspensions place $12\nu_9$ in the same
$\pi_{12}(S^9)\cong\mathbb Z/24$ in which $E^6[g]=0$.
With the original peripheral character, define
$\Psi_H(M)=\mu_\eta(\chi(M)[H])$. Then
\[
 \Psi_H(C_0,C_1,C'_2)=(12\nu_6,0,12\nu_6),\qquad
 \ker\Psi_H=\chi^{-1}(2\mathbb Z/12).
\]
Its coefficient is $12\chi(M)$ modulo twenty-four, so it retains
parity. It is not a lift of $\chi$ under reduction modulo twelve:
that reduction is zero. Both complete proofs are included in the
higher-rung and cone/sphere drafts. The nonzero class comes from the
specific retained-angle map, not from matching the number twenty-four
to a geometric dimension.

\noindent\textbf{5 September: the modulo-$840$ lattice and full Weyl quotient.}
The shared source \path{es_niemeier_shared.tex} under
\path{supporting_materials/workbench/research/} proves
\[
 R_{\mathrm{ES}}=A_5^4\perp D_4\subset N,\qquad
 [N:R_{\mathrm{ES}}]=72,\qquad \det N=1,
 \qquad \#\operatorname{Roots}(N)=144,
\]
\[
 O(N)=W(R_{\mathrm{ES}})\rtimes L(G),\qquad
 G\cong GL_2(\mathbb F_3)\cong Q_8\rtimes S_3.
\]
The abstract Weyl quotient is already recorded by
\href{https://arxiv.org/abs/1307.5793}{Cheng--Duncan--Harvey, Table~3}.
Here its code isomorphism and its full lattice lift are explicit in the
original fibre order. The theorems ``The exact full automorphism quotient
of $N$'' and ``The residue six-cycle in the full automorphism group''
prove that $\rho(S_{840})\subset W(R_{\mathrm{ES}})$,
$S_{840}=\langle289\rangle$, and that this six-cycle centralizes the
chosen $L(G)$. The $D_4$ discriminant map $G\to S_3$ has kernel $Q_8$.
The shortest dual classes $v,s,c$ give precisely the vector, even
half-spin and odd half-spin weights of the flag tangent modules described
by \href{https://www.uni-muenster.de/FB10/mjm/vol_6/mjm_vol_6_14.pdf}%
{Mare--Willems, Proposition~2.7 and equations (5.4)--(5.6)}.
For the fixed shell $R=3$, the residue-one primes $2521$ and $3361$
have raw counts $0$ and $6$, and unordered counts $0$ and $3$.
Thus the actual divisor counts do not descend to prime residues alone;
the finite-sample Fourier pushforward below retains those counts.

\noindent\textbf{5 September: the fifteen-selector test.}
At $p=1201$, the complete census maps by
$(h,r,t)\mapsto h\left(\begin{smallmatrix}r^2&rt\\rt&t^2\end{smallmatrix}\right)$
to fifteen distinct projective rays. Embedding these matrices in the
upper-left corner of $\operatorname{Herm}_3(\mathbb H)$ gives linear span
and Jordan closure exactly $\operatorname{Sym}_2(\mathbb R)$, of dimension
three. The full quaternionic Peirce corner has dimension six. For primes
$1201,4561,8761,9601$, all congruent to $361$ modulo $840$, the complete
oriented state counts are respectively $15,54,36,29$; the corresponding
unordered triple counts are $11,40,26,23$. The expanded divisor-channel
draft contains the full census, map, and closure proof.

\noindent\textbf{4 September arithmetic addition.}
For $p\equiv1\pmod4$, divisors $R\mid(p+1)/2$, $R\equiv3\pmod4$
give $x=(p+R)/4$, $k=(p+1)/R$ and
$4/p=1/x+1/(xk)+1/(pxk)$. Their number is exactly
$\frac12(\tau((p+1)/2)-\sum_{d\mid(p+1)/2}\chi_4(d))$; a finite-CRT
argument gives relative density one in every fixed reduced progression
of residue one modulo four, not every individual prime. The audit also gives
a rational inverse-bearing map from $ac-b^2=-\delta$ to $\Delta(q)=-1$:
with $D=-1-a^2-b^2+4a+4b+4ab$ and $K=1-2a-2b$ it sends
$A\mapsto(1+2K/D,2Ka/D,2Kb/D)$. For integral $a,b$, its exact image
denominator is $|D|/\gcd(|D|,|2K|)$. It is different from Pro's
orthogonal map, whose rational image is empty; neither map establishes
membership in the global period curve of the proposed threefold.

\noindent\textbf{4 September quaternionic continuation.}
The two Pro constructions define the same embedded thirteen-space:
$\{H:PHP^*\in\operatorname{Herm}_2(\mathbb C)\}
=\{\Theta(A,b):u^*b=0\}$, where $P$ has rows $(1,-1,0)$ and $(0,1,-1)$,
$u=(1,1,1)^{\mathsf T}$, and $\Theta(A,b)=A+B(b)\mathbf j$ with
$B(b)_{ij}=\sum_k\epsilon_{ijk}\bar b_k$. Their second spinors satisfy
$w=-\left(\begin{smallmatrix}0&1\\-1&0\end{smallmatrix}\right)\bar\chi$;
their common positive cone gives the exact boundary pair $(S^{12},S^8)$.
This geometric identification does not erase the dynamics of the omitted
complex coordinate. In the supplied scalar model, with
$s=\operatorname{tr}\Phi^2+2b^*b$ and
$V=m^2s/2+\Gamma s^2/4-\kappa(\det\Phi-b^*\Phi b)$,
\[
 \rho=u^*b,\qquad
 \ddot\rho=-(m^2+\Gamma s)\rho-\kappa u^*\Phi b.
\]
At zero initial velocities, $\Phi=\operatorname{diag}(1,-1,0)$,
$b=(1,-1,0)^{\mathsf T}$, one has $\rho=0$ but $\ddot\rho=-2\kappa$.
Nonzero motions can nevertheless remain in the slice without that reaction:
for each fixed complex $W\subset u^\perp$, keep $b\in W$,
$\Phi W\subset W$, and tangent initial velocities. The original equations
preserve this family globally; its configuration dimension is six for a
line $W$ and eight for $W=u^\perp$. The line case permits a nonzero marked
cross-coordinate $z$. In contrast, other initial data have vanishing
normal derivatives through order three but $\rho^{(4)}(0)=-12\kappa$.
Thus an instantaneous tangency check is weaker than the proved invariant
family. The proof supplement derives the constrained model's exact reaction force
$\kappa(uu^*/3)\Phi b$, retains the nonunitary Jordan metric
$D=\operatorname{diag}((PP^*)^{-1},3)$, and corrects the gauge kinetic sign
for anti-Hermitian generators. It makes no continuum mass-gap claim.

\noindent\textbf{Periodic continuation and the scalar vacuum.}
The fixed-subspace families are not exhaustive. For
$(m^2,\Gamma,\kappa)=(1,1,6)$ take
$b_0=(1,-1,0)^{\mathsf T}/4$ and
$\Phi_0=\frac14\left(\begin{smallmatrix}1&1&1\\1&1&1\\1&1&-2\end{smallmatrix}\right)$.
Let $O(t)$ be rotation about $(1,1,1)$ with angular frequency $\sqrt2$.
The explicit original-coordinate rotation is given in the proof supplement.
Then $b(t)=O(t)b_0$, $\Phi(t)=O(t)\Phi_0O(t)^*$ solve the full scalar
equations and have $u^*b=\Phi b=0$, $s=1$, and energy $7/4$.
The moving $b$ spans $u^\perp$, while $\Phi_0$ does not preserve that
plane. The conserved moment is $-\Omega/2$, $\Omega=\dot O O^{-1}$;
this nonzero matrix prevents the zero-connection lift from satisfying
Gauss' equation, even in the gauge where the scalar coordinates are static.
Separately, the supplied potential's nonzero minimum manifold is
$\{x(3P-I_3):P\in\mathbb{HP}^2\}$ below
$m^2=\kappa^2/(27\Gamma)$, with coexistence at equality. Here
$x=(\kappa+\operatorname{sgn}(\kappa)\sqrt{\kappa^2-24\Gamma m^2})/(12\Gamma)$.
Its full metric Hessian has eigenvalues $x(12\Gamma x-\kappa)$, $0$,
and $3\kappa x$ of multiplicities $1,8,5$. These are complete scalar
calculations, not a physical quantum mass-gap claim.

\overviewsection{1. The candidate, its fillings, and the original audit}

\noindent\textbf{The common higher-rung object.}
The nine-dimensional complex-Hermitian space has a marked $4+4+1$
decomposition; the two $4+4+4+1$ constructions are the same embedded
thirteen-slice of $\operatorname{Herm}_3(\mathbb H)$, and its Jordan
closure is the full $6+4+4+1$ space. A positive observation of that full
space, $H\mapsto\pi_{\mathbb C}(PHP^*)$, retains a four-dimensional
Lorentz block $G$ without imposing $\rho=u^*b=0$. Its exact permitted
pairs are $G\succeq0$, $|\rho|^2\leq\det G$: every such pair lifts.
Consequently a null block forces $\rho=0$, while timelike blocks admit
nonzero $\rho$. This inequality alone has no uniform positive threshold;
$G=\epsilon I_2$, $\rho=\epsilon/2$ has positive lifts for every
$\epsilon>0$. For the trace-free scalar variable $Q=\Theta(\Phi,b)$,
the separate positive-boundary map is $Q\mapsto Q-\lambda_{\min}(Q)I_3$,
and its observation is $P\Phi P^*-\lambda_{\min}(Q)PP^*$.
No nonzero trace-free matrix is itself positive. These exact maps
identify a coherent carrier and observation problem; the rank-six
integral toroidal extension supplies the separate period and filling
data needed for a successor to the original cusp construction.

\noindent\textbf{What now selects the observation, and what does not.}
Keep the exact rank-six periods $\Omega_3=(I_3,Z_3)$,
$Z_3=sI_3+C_3(t)$, $s=x+iy$, including the inherited correction $C_3$.
For real positive lattice metrics, the map
$W_3\mapsto H_3=\Omega_3W_3\Omega_3^*$ has a complete
twelve-dimensional positive fibre; the observed map $W_3\mapsto PH_3P^*$
has a seventeen-dimensional fibre. Every smooth positive $G(t)$ lifts
with the original monodromy. In particular
$G=\operatorname{diag}(\alpha_{\rm obs},\beta_{\rm obs}y^{-p})$
has an explicit positive equivariant lift and approaches a nonzero null
block. Its relative metric condition number grows as
$2\max\{2\alpha_{\rm obs}/3,3\lambda\}y^{p+2}/\beta_{\rm obs}$.
By contrast a bounded relative condition number excludes such rank loss.
The periods thus support the proposed observation but do not, without
a metric-selection law, determine its trajectory.

The full gauge equations give a second exact test. For the rotating
scalar family, temporal-gauge scalar compatibility forces
$\mathcal A_i=ic_i\Phi$; Gauss' residual is
$(2(x^2+d^2)\sum_i c_i^2/g^2+x^2+d^2+B^2)\Omega\ne0$.
Arbitrary temporal connection is retained in a separate exact inertia
equation; no general impossibility is inferred from temporal gauge.
There is a constructive full-gauge sector with
$\Phi=x(t)\operatorname{diag}(-2,1,1)$,
$b=\beta(t)e_1$, $\mathcal A_\mu=0$:
\[
 \ddot x=-ax-\kappa x^2+\kappa\beta^2/3,\qquad
 \ddot\beta=-(a-2\kappa x)\beta,\qquad
 a=m^2+\Gamma(6x^2+2\beta^2).
\]
Every real initial condition gives a global solution, including the
exact scalar minimum $x=(3+\sqrt3)/12$, $\beta=3x$ at $(m^2,\Gamma,\kappa)=(1,1,6)$.
For a quaternionic positive rank-one $H=\Theta(A,b)$, the gauge-invariant
average of the marked determinant is
\[
 \int_{SU(3)}\det(PgAg^*P^*)\,d\mu(g)
 =V_\mu\operatorname{tr}\operatorname{adj}(A)=V_\mu\|b\|^2.
\]
An explicit family of equal-potential minima has
$\|b\|^2=9\chi^2\sin^2(2\theta)/4$, $0\leq\theta\leq\pi/4$,
$\chi=(3+\sqrt3)/6$. The fixed-mark determinant can change from timelike
to null inside one gauge orbit, whereas this average is invariant and
distinguishes those minimum orbits. None of these determinants is
identified here with a quantum Hamiltonian mass.

\noindent\textbf{Exact dynamics and the invariant vector mass.}
The later calculation derives an actual linearized gauge-vector mass
from the same invariant. Retain $m^2\geq0$, $\Gamma>0$, $\kappa>0$,
$m^2\leq\kappa^2/(27\Gamma)$ and
$\chi=(\kappa+\sqrt{\kappa^2-24\Gamma m^2})/(12\Gamma)$.
For $q=\cos(2\theta)$, $0\leq\theta\leq\pi/4$, use the minimum
\[
 \Phi_\theta=\chi\operatorname{diag}
       ((1+3q)/2,(1-3q)/2,-1),\qquad
 b_\theta=-\tfrac32\chi\sin(2\theta)e_3 .
\]
With the retained anti-Hermitian kinetic inner product $-\operatorname{tr}(ZW)$,
all eight inertia eigenvalues, including multiplicities, are
\[
 \chi^2\{0,\;3(1-q^2),\;
       9q^2\ [2],\;\tfrac92(1+q)\ [2],\;\tfrac92(1-q)\ [2]\}.
\]
Their transverse squared masses are $g^2/2$ times these values: direct
substitution into all linearized scalar equations and Gauss gives
$\omega^2=|k|^2+g^2\lambda/2$. The exact nonlinear phase current
$j_\nu=-i(b^*D_\nu b-(D_\nu b)^*b)$ is gauge invariant and conserved.
For a transverse fluctuation along $i\operatorname{diag}(1,1,-2)$,
it is $-4\|b_\theta\|^2 a_\nu$ at first order. This detects the branch
\[
 \mu_{\rm d}^2=\frac{2g^2}{3}\|b_\theta\|^2
 =\frac{2g^2}{3V_\mu}
       \int_{SU(3)}\det(PgA_\theta g^*P^*)\,d\mu(g).
\]
Here $A_\theta=\Phi_\theta+\chi I_3$ is the complex part of the
positive rank-one quaternionic matrix. The current coefficient vanishes
at $b_\theta=0$; it is not a nonzero linear interpolator at that endpoint.
One transverse zero branch persists in the interior, with kernel
dimensions four and three at the two endpoints.

There are also exact nonlinear full-gauge continuations with zero
connection. In the original scalar coordinates put
\[
 \Phi=\operatorname{diag}(h+k,h-k,-2h),\qquad b=\beta e_3.
\]
All gauge currents vanish for any real profiles. The complete remaining
equations are
\[
 \Box h=-ah-\kappa h^2+\kappa(k^2+\beta^2)/3,\quad
 \Box k=-(a-2\kappa h)k,\quad
 \Box\beta=-(a-2\kappa h)\beta,\quad
 a=m^2+\Gamma(6h^2+2k^2+2\beta^2).
\]
For $0<m^2<\kappa^2/(27\Gamma)$ set $h_0=\chi/2$ and choose
$h_0<h<\kappa/(6\Gamma)$. The exact solutions are
\[
 a=\frac{m^2+12\Gamma h^2}{1-6\Gamma h/\kappa},\quad
 R^2=\frac{3ah}{\kappa}+3h^2,\quad
 \omega^2=a-2\kappa h>0,\quad
 k(t)=R\cos(\omega t),\quad\beta(t)=R\sin(\omega t).
\]
As $h\downarrow h_0$, their frequency tends to zero and their
energy density satisfies $E-V_0=9h_0^2\omega^2+O(\omega^4)$.
The amplitude tends to $3h_0$, not zero: a finite gauge-invariant
excursion becomes arbitrarily slow. At fixed minimum amplitude,
$\theta=f(t-x_1)$ instead gives an exact nonlinear travelling wave
for any smooth $f$. On a rectangular torus and $L_1$-periodic $f$,
its excess energy is $18\chi^2L_2L_3\int_0^{L_1}(f')^2$.
These energy statements retain finite spatial volume. They neither
quantize the model nor decide its interacting quantum spectrum.

\noindent\textbf{The full homogeneous quantum model now has a proved gap.}
The separate draft \path{s6_homogeneous_quantum_confinement.tex}
quantizes all three $\mathfrak{su}(3)$ spatial gauge matrices, all eight
adjoint-scalar coordinates, and all six triplet coordinates. With
$T_A=-\Delta_{\mathfrak{su}(3)^3}$ and
$T_E=-\Delta_{\operatorname{tr}(d\Phi^2)+2db^*db}$, its Hamiltonian is
\[
 H=\frac{\hbar^2g^2}{4\mathcal V_3}T_A+
       \frac{\hbar^2}{2\mathcal V_3}T_E+\mathcal V_3U,
\]
\[
 U=\frac1{g^2}\sum_{i<j}\|[A_i,A_j]\|^2
   +\frac12\sum_i\operatorname{tr}[A_i,\Phi]^2
   +\sum_i\|A_i b\|^2+V(\Phi,b).
\]
Nothing is restricted to the minimum curve. The exact root sum
$\operatorname{tr}_{\mathfrak{su}(3)}\sqrt{-\operatorname{ad}_X^2}
=4(\lambda_{\max}(-iX)-\lambda_{\min}(-iX))\geq2\sqrt6\|X\|$
gives, for every $0<\eta<1$,
\[
 H\geq\frac{\eta\hbar^2g^2}{4\mathcal V_3}T_A+
 \frac{\hbar^2}{2\mathcal V_3}T_E+
 \hbar\sqrt{6(1-\eta)}\sum_i\|A_i\|
 +\frac{\mathcal V_3\Gamma s^2}{8}
 -\frac{\mathcal V_3\kappa^4}{54\Gamma^3}.
\]
This quadratic-form bound comes from transverse quantum oscillators,
not an added classical potential. It proves compact resolvent despite
commuting classical valleys. The derived Gauss generators impose
spatially constant $SU(3)$ invariance, and the strictly positive heat
kernel proves a unique invariant ground state. Thus
$E^{\rm phys}_1-E_0>0$ for fixed positive
$\hbar,g^2,\mathcal V_3,\Gamma$, $m^2\geq0$ and real $\kappa$.
The bosonic matrix mechanism is classical (Simon, 1983, Section~7);
the calculation applies it with full constants to this coupled action.
No numerical or volume-uniform gap follows. The classical embedding
into torus fields is proved, but it is not an automatic quantum
subsector: the Fourier pair
$\Phi=(q+y\cos(2\pi x_1/L_1))\operatorname{diag}(1,-1,0)/\sqrt2$
has the mixed potential $3\Gamma\mathcal V_3q^2y^2/4$. It prevents
all homogeneous wavefunctions tensored with one fixed normalizable
discarded-mode state from being a reducing subspace. The continuum
comparison must retain this coupling.

\noindent\textbf{The period geometry now has a precise spectral meaning.}
The original $W$ is the inverse metric in lattice coordinates.
Hermitian compatibility gives $h_3=2H_3^{-1}$ and a quotient metric
$h_2=2G^{-1}$; the original map $P\oplus Q_{\rm latt}$ intertwines
the Laplacians exactly on pulled-back characters. For the native
Euclidean source, $H_3=2I_3$ and $G=2PP^*$ are constant, whereas
\[
 \lambda_{m,n}=4\pi^2\left(\|m\|^2+
       \|(yI_3+\operatorname{Im}C_3)^{-\mathsf T}
          (n-(xI_3+\operatorname{Re}C_3)^{\mathsf T}m)\|^2\right).
\]
The source and induced quotient first positive eigenvalues are
$4\pi^2/y^2$ and $8\pi^2/y^2$ to leading order, with explicit two-sided
bounds retaining the full correction $C_3$. Thus constant observed
cometric does not prevent the native periods from expanding and their
positive Laplace gaps from tending to zero.
The full proofs are in the quaternionic continuation draft, including
the separate exclusions of all time-dependent real gauge lifts and
constant-spatial complex lifts of the earlier literal rotating orbit.

\noindent\textbf{A controlled quantum correspondence retaining the Fourier interaction.}
In the exact truncation $\Phi=(q+y\cos(2\pi x_1/L_1))
\operatorname{diag}(1,-1,0)/\sqrt2$, $A_i=b=0$, the actual fibre
operator and its ground state are
\[
 h(q)=-\frac{\hbar^2}{\mathcal V_3}\partial_y^2
       +\frac{\mathcal V_3\omega(q)^2}{4}y^2
       +\frac{3\Gamma\mathcal V_3}{32}y^4,
 \qquad h(q)\Omega_q=e(q)\Omega_q,
\]
where $\omega(q)^2=\omega_*^2+3\Gamma q^2$,
$\omega_*^2=m^2+(2\pi/L_1)^2$. The strictly positive ground state,
normalized in $\sqrt{\mathcal V_3/2}\,dy$, is not Gaussian.
The isometry $Jf=f\Omega_q$ gives the exact compressed Hamiltonian
\[
 H_{\rm eff}=-\frac{\hbar^2}{2\mathcal V_3}\partial_q^2
 +\mathcal V_3\left(\frac{m^2q^2}{2}+\frac{\Gamma q^4}{4}\right)
 +e(q)+\frac{\hbar^2}{2\mathcal V_3}\|\Omega'_q\|^2.
\]
Its Berry connection vanishes in the positive real phase, but its last
potential term is strictly positive for $q\ne0$. With
$Q_q=I-|\Omega_q\rangle\langle\Omega_q|$, the full residual is
\[
 H_2Jf-JH_{\rm eff}f
 =-\frac{\hbar^2}{2\mathcal V_3}
       (2f'\Omega'_q+fQ_q\Omega''_q).
\]
This map is not an invariant subsector. Nevertheless, for
$\mathcal V_3\geq9\Gamma\hbar/(16\omega_*^3)$, a normalized
initial state of effective energy expectation at most $E$ satisfies
\begin{align*}
 \|e^{-itH_2/\hbar}Jf-Je^{-itH_{\rm eff}/\hbar}f\|
 &\leq |t|\left[
 \frac{\sqrt{30\Gamma E}}{\omega_*\sqrt{\mathcal V_3}}
 +\frac{9\sqrt5\,\Gamma\hbar}{\mathcal V_3\omega_*^2}\right.\\[-1mm]
 &\hspace{22mm}\left.
 +\frac{96}{5^{3/4}}
      \sqrt{\frac{\Gamma\hbar}{\mathcal V_3\omega_*}}\right].
\end{align*}
The quartic fibre spectrum, derivative estimates and operator domains
are proved in the quantum draft. The bound improves at fixed
$\hbar,\Gamma,L_1,m^2,E,t$ as volume grows; it is not uniform when
the original Fourier frequency tends to zero and does not assert a
continuum spectral limit. This gives an explicit relation where the
fixed-vacuum product failed, rather than discarding that interaction.

\noindent\textbf{The remaining rotating-gauge equations are now explicit.}
For that literal scalar orbit, a genuinely complex class with temporal
and one spatial matrix exchanging the fixed matter line with its
orthogonal complement, and the other spatial matrix preserving that
line, is excluded by the full Gauss and spatial equations. Here the
unchanged constant-matter data are
\[
 H=\begin{pmatrix}0&0&0\\0&-1/2&1/(2\sqrt2)\\
       0&1/(2\sqrt2)&1/2\end{pmatrix},\quad
 q=\frac{(1,0,0)^{\mathsf T}}{2\sqrt2},\quad
 A=\sqrt2\begin{pmatrix}0&-1&0\\1&0&0\\0&0&0\end{pmatrix}.
\]
The exclusion
also covers every constant spatial rotation of that class. For arbitrary
$Z_1,Z_2\in\mathfrak{su}(3)$ in the two-spatial-rotation ansatz, retain
$\mathcal IX=([H,[H,X]]+qq^*X+Xqq^*)_{\rm tf}$ and
$\mathcal KX=-\sum_j[Z_j,[Z_j,X]]$. Every possible solution has
$[Z_1,Z_2]\ne0$, which makes $L=g^2\mathcal I+2\mathcal K$ positive
definite. The exact Gauss and scalar equations then give
\[
 T=4L^{-1}[Z_1,Z_2],\qquad C=\omega T,\qquad
 \omega^2=\tau:=\frac{1/4+\|Z_1q\|^2+\|Z_2q\|^2}{\|Tq\|^2}.
\]
The numerator retains the original $2q^*q=1/4$. This does not replace
the scalar equation by its contraction: all entries of
\[
 \tau[T,[T,H]]-\sum_j[Z_j,[Z_j,H]]-[A,[A,H]],\qquad
 (\tau T^2-Z_1^2-Z_2^2)q+2q
\]
must vanish, as must the two full spatial residuals printed in the
quaternionic draft. Their common zero set has an explicit inverse
$\omega=\pm\sqrt\tau$, $C=\omega T$ recovering the original fields.
This removes the temporal matrix and frequency from the search without
discarding colour coordinates or a field equation. It is not yet a
solution or an exclusion of that unrestricted zero set.
The full scalar equations also force the original stabilizer
coefficients $s_X=-i\operatorname{tr}(HX)$ to obey
$s_C^2<s_{Z_1}^2+s_{Z_2}^2$ unless all three vanish.
This strict signature constraint is proved for every complex triple,
not just the excluded block class. A further exact two-by-two stress
equation, including its sum-of-squares discriminant, retains both
spatial equations as necessary tests of the remaining zero set.


\paragraph{5 September: finite coupling fibres and a proper spatial exceptional set.}
Keep every entry of the preceding $Z_1,Z_2$ and put
$t=g^2$, $J=iH$, $W=J^\perp$, and
$\Pi X=X-\langle J,X\rangle J/(3/4)$. Write
$\mathscr B=[Z_1,Z_2]$, $\eta=\langle J,\mathcal KJ\rangle$,
$w=\Pi\mathcal KJ$, and $\beta=\langle J,\mathscr B\rangle$.
For a noncommuting pair, $\eta>0$, and the exact seven-dimensional
pencil and forcing vector are
\[
 \mathscr M(t)=t\mathcal I|_W+
  2\left(\Pi\mathcal K|_W-\frac{w\langle w,\cdot\rangle}{\eta}\right),
 \qquad f_7=\Pi\mathscr B-\frac\beta\eta w.
\]
Define $D_7=\det\mathscr M$, $V_7=4\operatorname{adj}(\mathscr M)f_7$,
and retain the full temporal numerator
\[
 N_7=V_7+\frac{2\beta D_7-\langle w,V_7\rangle}{\eta}J,
 \quad Q_7=\|V_7q\|^2,
 \quad\alpha=\tfrac14+\sum_j\|Z_jq\|^2.
\]
Then $T=N_7/D_7$, $D_7(t)>0$ for $t>0$, and the determinant has
degree seven with leading coefficient $3/128$.
Clearing only the retained positive frequency denominator in all four
full residuals gives $(R_H,R_q,R_1,R_2)=(P_H,P_q,P_1,P_2)/Q_7$.
Their original Hermitian/vector/anti-Hermitian codomains give the
nonnegative polynomial
\[
 p_Z(t)=\operatorname{tr}(P_H^2)+P_q^*P_q-
                 \operatorname{tr}(P_1^2)-\operatorname{tr}(P_2^2).
\]
For $Q_7\not\equiv0$ its degree is exactly twenty-eight, with
\[
 [t^{28}]p_Z=\alpha^2(3/128)^4
       \|\mathcal L_{2\beta/\eta}^2(Z_1,Z_2)\|^2>0,\qquad
 \mathcal L_a(X_1,X_2)=(X_1+a[J,X_2],X_2-a[J,X_1]).
\]
Every real root has even multiplicity. Thus at most fourteen
positive couplings satisfy the original system over a fixed pair;
the exact permitted roots retain $Q_7(t)>0$ and the explicit bounded
interval proved in the quaternionic draft. Each such root has precisely
the two reconstructions
\[
 \omega=\varepsilon\sqrt\alpha D_7/\sqrt{Q_7},\qquad
 C=\varepsilon\sqrt\alpha N_7/\sqrt{Q_7},\qquad\varepsilon=\pm1.
\]
The identically empty domain is itself decided by three vector tests:
$f_7q=(\mathcal S_7f_7)q=(\mathcal S_7^2f_7)q=0$, where
$\mathcal S_7$ is the second summand of $\mathscr M/2$ after
removing $t\mathcal I|_W/2$.

There is also a global restriction on the sixteen-coordinate spatial
projection. In the fixed diagonal frame $U$ already used for $H$, the
two polynomials
$F_\xi=\eta^{16}(U^*P_HU)_{11}$ and
$G_\xi=\eta^{16}(U^*P_HU)_{22}$ have degree at most twelve in $t$
and polynomial coefficients in all spatial coordinates $\xi$.
The $24\times24$ matrix of $(A,B)\mapsto AF_\xi+BG_\xi$, with
$\deg A,\deg B<12$, has determinant $\mathscr R(\xi)$.
Every full solution satisfies $\mathscr R(\xi)=0$.
For the literal full-coordinate pair
$(a,h,v_1,v_2,z)=(1,2,1+i,-1+2i,2-i)$ and
$(-2,1,2+3i,1-i,-1+2i)$, an exact rational Bezout identity shows
this determinant is nonzero. Independently, reduction through
$\sqrt6\mapsto39$ modulo $101$ gives determinant $60\ne0$ after
the displayed $\eta^{16}$ multiplication. The complete determinant
definition, inverse coordinate maps, and certificates are in the draft.
Consequently solutions can project only into a closed, nowhere-dense,
measure-zero algebraic set. Existence on that set requires further analysis.

\paragraph{5 September: an excluded subspace and an exact error bound.}
This exceptional set contains a whole eight-dimensional linear subspace:
put $P=qq^*/(q^*q)$ and let both $Z_1,Z_2$ exchange $\mathbb Cq$
and $q^\perp$, so $PZ_jP=(I-P)Z_j(I-P)=0$.
The Gauss equation forces $C$ to preserve these two summands.
Two explicitly invertible simultaneous unitary changes then give
\[
 C=ie\operatorname{diag}(1,0,-1),\quad
 H=\begin{pmatrix}0&0&0\\0&0&r\\0&r&0\end{pmatrix},\quad
 q=e_1/(2\sqrt2),\quad r=\sqrt6/4,
\]
with neither the sign of $e$ nor the sign of $\omega$ selected.
Every entry of the spatial vectors is retained in
$Z_1=X(v)$, $Z_2=X(w)$, $X(v)=\left(\begin{smallmatrix}0&-v^*\\v&0\end{smallmatrix}\right)$.
Write $M=vv^*+ww^*$, $u=M_{11}$, $v_0=M_{22}$, $m=M_{12}$ and
$d_M=\det M$. Both scalar equations give
$e^2=2+\operatorname{tr}M$ and $\operatorname{Re}m=\delta=1/(2r)>0$.
With $N=i(vw^*-wv^*)$ and $b=\operatorname{Im}(v^*w)$, full Gauss gives
\[
 b=-\frac e{2\omega}(u+2v_0+t/24),\quad
 N_{12}=-\frac{3e}{2\omega}m,\quad
 N_{22}=-\frac e\omega(2v_0+5t/24),\qquad t=g^2>0.
\]
The complete spatial residual vectors satisfy
\[
 s_1v^*+s_2w^*=DM+d_MI_2+(4bI_2-2\omega eQ)N,\quad
 Q=\operatorname{diag}(1,2),\quad D=(t/4-\omega^2)I_2-e^2Q^2.
\]
If both spatial equations held, the $(2,1)$ and $(2,2)$ entries
would imply
\[
 0=d_M+\frac{e^2}{\omega^2}
       (u+2v_0+t/24+2\omega^2)(v_0+5t/12)>0,
\]
a contradiction for every permitted coupling and both frequency signs.
Without setting the spatial residuals to zero, their full matrix versions
$R_j=X(s_j)$ satisfy, on the same scalar--Gauss zero set,
\[
 \sqrt{\|R_1\|_{\mathrm{HS}}^2+\|R_2\|_{\mathrm{HS}}^2}
 >\frac{5\sqrt2\,e^2t}
 {6\sqrt{\operatorname{tr}M(1+v_0^2/|m|^2)}}.
\]
For $\operatorname{tr}M\leq L$ and $t\geq t_0>0$ this exceeds
$5\sqrt2\,t_0/(3\sqrt{L(1+L^2/\delta^2)})$.
The full quaternionic draft proves a denominator-free identity in the
nonzero Gauss and spatial residuals, then derives this bound; it is not
a numerical tolerance chosen for a search. The constrained domain is
now completely parameterized and is nonempty at every original $t>0$,
as the following continuation proves.

The subspace lies in $\{\mathscr R=0\}$ because the four-dimensional
cross block of the Gauss pencil contributes a degree-four determinant
$D_4$ with leading coefficient $1/16$, and both selected diagonal scalar
polynomials contain $D_4^2$. This inclusion is proved through an explicit
kernel of the fixed Sylvester map, including degree drops and the zero
pair. It does not imply these common factors have admissible positive roots.
The complete proof excludes this subspace, not the remaining exceptional
set or arbitrary time-dependent complex gauge lifts.

\paragraph{5 September: every-coupling scalar--Gauss classification.}
The entire classified scalar--Gauss domain is in bijection with
$\mathcal D\times\{1,-1\}^2\times U(1)^2$, where
\[
 \mathcal D=\{(u,v_0,y)\in\mathbb R^3:
 u>0,\ d=uv_0-\delta^2-y^2>0,\ c=4v_0-5u>0\},\quad
 \delta=\sqrt{2/3}.
\]
Retain $m=\delta+iy$ and set
\[
 \begin{gathered}
 M=\begin{pmatrix}u&m\\\bar m&v_0\end{pmatrix},\quad t=72d/c,\qquad
 T=\frac t4E-\frac12(QM+MQ),\\
 Q=\operatorname{diag}(1,2),\qquad E=\operatorname{diag}(2/3,-5/6).
 \end{gathered}
\]
\[
 \begin{gathered}
 e=\epsilon\sqrt{2+u+v_0},\qquad
 k=\sigma\sqrt{\frac d{-\det T}},\qquad
 \omega=e/k,\\
 N=kT,\qquad \epsilon,\sigma\in\{1,-1\}.
 \end{gathered}
\]
The full proof establishes $\det T<0$, and $A_\pm=M\pm N$ are
positive semidefinite of rank one with $h_\pm=(A_\pm)_{11}>0$. Thus
\[
 \begin{gathered}
 a_\pm=\begin{pmatrix}\sqrt{h_\pm}\\(A_\pm)_{21}/\sqrt{h_\pm}\end{pmatrix},\qquad
 p_\pm\in U(1),\\
 v=\frac{p_+a_++p_-a_-}2,\quad
 w=\frac{i(p_+a_+-p_-a_-)}2
 \end{gathered}
\]
are defined throughout $\mathcal D$, and
$p_\pm=(v\mp iw)_1/\sqrt{h_\pm}$ recovers the phases exactly.
Their full scalar residuals vanish because
$e^2=2+\operatorname{tr}M$, $\operatorname{Re}m=\delta$.
The complete lower Gauss block is
\[
 etE-2e(QM+MQ)-4\omega N
 =etE-2e(QM+MQ)-4e\left(\frac t4E-\frac12(QM+MQ)\right)=0;
\]
the remaining upper entry is minus its trace, and all off-block
entries are zero. These are full residual identities.

Gauss and scalars force every literal temporal matrix into the displayed
classified form by the retained inverse frame $Y_0=B(\zeta)Y'B(\zeta)^*$.
Writing $B(\zeta)=U\operatorname{diag}(1,W(\zeta))$ in the draft's
fixed first frame,
\[
 W(\zeta)=\frac1{\sqrt2}
 \begin{pmatrix}-\zeta^{-1}&-\zeta^{-1}\\\zeta&-\zeta\end{pmatrix},
 \qquad |\zeta|=1,
\]
retains both square roots of the temporal phase. The complete literal
domain is therefore in bijection with
\[
 \bigl(\mathcal D\times\{1,-1\}^2\times U(1)^2\times U(1)\bigr)
 /\bigl((p_+,p_-,\zeta)\sim(-p_+,-p_-,-\zeta)\bigr).
\]
The proof computes every inverse coordinate and shows that each raw
fibre consists of precisely these two points; no phase is discarded.
The smooth coordinates
$(\eta,\alpha,\beta)=(\zeta^2,p_+/\zeta,p_-/\zeta)$ give
\[
 \mathcal S_{\rm sg}\cong\{1,-1\}^2\times(0,\infty)
                         \times\mathbb R^2\times\mathbb T^3.
\]
For $(\tau,y,s)\in(0,\infty)\times\mathbb R^2$ the exact coupling
coordinates are
\[
 v_0=\frac{\sqrt5}{2}\sqrt{\delta^2+y^2}+e^s,\qquad
 u=\frac{4\tau v_0+72(\delta^2+y^2)}{72v_0+5\tau},\qquad t=\tau=g^2.
\]
The source \path{scalar_gauss_smooth_shared.tex} proves smooth inverse
charts in the literal matrix entries and embeddedness. There are exactly
four six-dimensional literal components; every fixed positive coupling
has four embedded five-dimensional components
$\mathbb R^2\times\mathbb T^3$. The positive coordinate is the unchanged
coupling, and no physical gauge quotient or spatial equation is imposed.
For every prescribed original $\tau=g^2>0$, take
\[
 \rho=\frac{\tau+\sqrt{\tau^2+4608\delta^2}}{96\delta}
       >1/\sqrt2,\qquad
 M_\rho=\delta\begin{pmatrix}\rho&1\\1&2\rho\end{pmatrix}.
\]
Here $d=\delta^2(2\rho^2-1)>0$, $c=3\delta\rho>0$, and
$48\delta\rho^2-\tau\rho-24\delta=0$ gives $t=72d/c=\tau$.
The independent signs retain both temporal and original frequency signs.
This proves a nonempty domain for the strict spatial-error certificate;
the two spatial residuals cannot vanish simultaneously throughout the class. Existence
on the remaining unrestricted exceptional set stays open. These classical
residual statements do not establish a quantum gap.
The complete proofs and independent exact replay are retained under
\path{supporting_materials/workbench/evidence/scalar_gauss_domain_20260905/}.



The construction starts from
\[
 B^\circ=\mathbb P^1\setminus\{p_0,p_1,p_2\},\qquad
 \Delta\simeq(\mathbb Z/3)*(\mathbb Z/4),\qquad g_0=(g_1g_2)^{-1},
\]
and the marked two-tori
\[
 \Pi=[Z\mid I_2],\quad
 Z=\begin{pmatrix}6\mu&\tau\\ \beta&\mu\end{pmatrix},\quad
 D=\operatorname{Im}\beta-
 \frac{6(\operatorname{Im}\mu)^2}{\operatorname{Im}\tau}<0,\quad
 \det_{\mathbb R}\Pi=\operatorname{Im}\tau\,D<0.
\]
The elliptic orientations are simultaneous.  With $s_j$ the Cayley
coordinates and
$S_1(w)=(w-1)/w$, $S(w)=-1/w$, $P(w)=w-1$,
\[
 \begin{gathered}
 s_1\circ g_1=e^{-2\pi i/3}s_1,\qquad
 s_2\circ g_2=-is_2,\qquad g_0\ \text{parabolic},\\
 \tau\circ g=\varphi(g)\circ\tau,\qquad
 (\varphi(g_1),\varphi(g_2),\varphi(g_0))=(S_1,S,P),\qquad S_1SP=1.
 \end{gathered}
\]
This normalization is forced at the cusp.  On a $g_0$-stable component
$U_0$, choose the $j$-cusp component $K$ containing $\tau(U_0)$ and put
\[
 P(K)=K,\qquad
 \tau'=\gamma_0^{-1}\tau|_{U_0},\qquad
 P'=\gamma_0^{-1}P\gamma_0,\qquad \tau'\circ g_0=P'\circ\tau'.
\]
If $P'=1$, then $\exp(2\pi i\bar\tau')$ has a single-valued logarithm and
winding zero, whereas the pole of $j(\tau)$ forces a zero and positive
winding.  Hence $P'$ is nontrivial parabolic.  Returning to
$\tau(z_1)=\rho$, $\varphi(g_1)=S_1$, the order-two condition gives
\[
 0=\operatorname{tr}(\widetilde S_1^{-1}\widetilde P^{-1})
 =1+k(p^2-pr+r^2),\qquad k=-1,\qquad p^2-pr+r^2=1,
\]
so conjugation by a power of $S_1$ gives the displayed triple.  Only then
is $q:=e^{2\pi i\tau}=t_cu_1(t_c)$, $u_1(0)\ne0$, redefined.  In the
resulting cusp coordinate $e^{2\pi is}=t_c$.

For $M_g=\rho_V(g)^{\mathsf t}$ and $A_g=M_g^{-1}$, the exact deck identity
and determinant signs are
\[
 \Pi(gz)=R_g(z)\Pi(z)M_g,\quad
 R_{g_1}=\begin{pmatrix}-1/\tau&0\\(1-\mu)/\tau&1\end{pmatrix},\quad
 R_{g_2}=\begin{pmatrix}1/\tau&0\\-\mu/\tau&1\end{pmatrix},
\]
\[
 \det R_{g_1}=-\tau^{-1},\qquad
 \det R_{g_2}=\tau^{-1},\qquad \det R_{g_0}=1.
\]
At the cusp,
\[
 Z=sB_0+C(t_c),\quad B_0=\begin{pmatrix}0&1\\-1&0\end{pmatrix},\quad
 \Psi_{\bar\lambda}(x,t_c)=
 \bigl(e^{2\pi iC(t_c)\bar\lambda}t_c^{B_0\bar\lambda}x,t_c\bigr),\quad
 N_0=Y_\epsilon/\bar\Lambda.
\]
The reduced irreducible fibre $W$ has local equations
$z_0$, $z_0z_1$, $z_0z_1z_2$ and normalization $\mathrm{dP}_6$ with
opposite hexagon sides identified.  The invariant toric form gives
\[
 \Omega=\frac{dx_1\wedge dx_2\wedge dt_c}{x_1x_2}
 =\pm dz_0\wedge dz_1\wedge dz_2,\qquad
 \Psi_{\bar\lambda}^*\Omega=\Omega,\qquad
 \omega_X|_{N_0}\simeq\omega_{N_0}\simeq\mathcal O_{N_0}.
\]

The opposite-edge identifications are literal holomorphic maps. For
$C_v=E_0\cap E_v$, put
$\iota_v=\Psi_{B_0v}|_{C_v}:C_v\to C_{-v}$, and use the same
primitive character $w=\chi^{(a,0)}$, $av=0$, on both dense edge orbits.
Then $\iota_v^*w_{-v}=e^{2\pi i\,aC(0)B_0v}w_v$, with
\[
 \begin{array}{c|c|c}
 v&a&\iota_v^*w_{-v}/w_v\\\hline
 e_1&(0,1)&e^{-2\pi i\mu_0}\\
 e_2&(1,0)&e^{12\pi i\mu_0}\\
 e_1-e_2&(1,1)&e^{-2\pi i(7\mu_0+b_0)}
 \end{array}
\]
where $\mu_0=\mu(0)$ and $b_0=b(0)$. The reverse maps have reciprocal
multipliers; these restrictions extend to the whole rational curves.
The full matrix $C(0)$ is retained in this calculation.

For $(m_1,m_2)=(3,4)$ and $(v_1,v_2)=(\varepsilon,-\varepsilon')$, the
finite lifts are
\[
 \widetilde g_j^b(z,\zeta)
 =(g_jz,R_{g_j}(z)\zeta+\Pi(g_jz)b),\qquad
 Q_{j,\log}=\widetilde g_j^{\,v_j/m_j},\qquad
 Q_{j,\mathrm{lin}}=\widetilde g_j^{\,0},
\]
\[
 Q_{j,\log}\tau_\lambda Q_{j,\log}^{-1}=\tau_{A_j\lambda},\qquad
 (Q_{j,\log})^{m_j}=\tau_{v_j},\qquad
 (Q_{j,\mathrm{lin}})^{m_j}=1.
\]
The invariant covector satisfies
$\gamma A_j=\gamma$,
$\gamma(kv_1/3)=k/3$ and $\gamma(kv_2/4)=-k/4$ for $0<k<m_j$;
hence every nonidentity power of the descended $q_{j,\log}$ is fixed-point
free.  Thus
\[
 N_j=J'_j/\langle q_{j,\log}\rangle,\qquad f_j^*(p_j)=m_jS_j,
\]
and the punctured conjugacy is
$h_jq_{j,\log}h_j^{-1}=q_{j,\mathrm{lin}}$ on $J_j^{\prime *}$.
These maps glue a compact smooth threefold
$f:X=N_0\cup J\cup N_1\cup N_2\to\mathbb P^1$.
The integer
\[
 p=12\ell_0-4\ell_1-3\ell_2,\qquad
 (\ell_0,\ell_1,\ell_2)=(0,1,-1)\Longrightarrow p=-1
\]
controls the topology.  The manuscript claimed
$X\cong_{\rm diff}S^6$, $a(X)=1$, and
$\langle c_3(T_X),[X]\rangle=2$; the first audits verified the integer
algebra but had not yet identified its collapse and matrices with the
actual specialization and oriented collar maps.

\clearpage
\overviewsection{2. Two integral topology certificates and recognition of $S^6$}

Let $F$ be a smooth fibre,
$V=H^1(F;\mathbb Z)=\mathbb Z\langle\gamma,u,w,\delta\rangle$, and
$q=uw+6\gamma\delta$.  Integral nearby cycles prove
\[
 \operatorname{sp}^r:H^r(W;\mathbb Z)\longrightarrow H^r(F;\mathbb Z),
 \qquad \operatorname{im}\operatorname{sp}^r
 =H^r(F;\mathbb Z)^{T_0}\quad(0\le r\le4),
\]
with saturated image.  In
$0\to R^rf_*\mathbb Z\to j_*L^r\to Q^r\to0$,
\[
 \begin{aligned}
 Q^1&=(\mathbb Z/3)_{p_1}\oplus(\mathbb Z/4)_{p_2},&
 Q^2&=Q^3=(\mathbb Z/2)_{p_2},&
 Q^4&=(\mathbb Z/3)_{p_1}\oplus(\mathbb Z/4)_{p_2},
 \end{aligned}
\]
and therefore
\[
 H^0(R^1f_*\mathbb Z)=\mathbb Z(12\gamma),\quad
 H^0(R^2f_*\mathbb Z)=\mathbb Z(2q),\quad
 H^0(R^3f_*\mathbb Z)=\mathbb Z(2\gamma uw).
\]
Canonical truncation identifies the disc--pair-of-pants discrepancy with
Leray $d_2$:
\[
 d_2(12\gamma)=-p\omega,\qquad
 d_2(2q)=-p\omega[\delta],\qquad
 d_2(2\gamma uw)=-p\omega[\gamma\delta].
\]
For $p=-1$ these maps are primitive, so
$H^1(X;\mathbb Z)=H^2(X;\mathbb Z)=H^3(X;\mathbb Z)=0$ and $H_2(X)=0$.
The actual based overlap maps give
\[
 \pi_1(X)=
 \langle c,x,y\mid c\ {\rm central},\ xy=1,\ x^3=c,\ y^4=c^{-1}\rangle=1.
\]
Consequently
\[
 H_k(X;\mathbb Z)=
 \begin{cases}\mathbb Z,&k=0,6,\\0,&1\le k\le5,\end{cases}
 \qquad X\simeq S^6,\qquad
 X\cong_{\rm diff}S^6,\qquad
 \langle c_3(T_X),[X]\rangle=\chi(X)=2,
\]
where the last diffeomorphism uses $\Theta_6=0$.

The independent cellular certificate starts from the equivariant honeycomb
collapse:
\[
 0\longrightarrow\operatorname{im}(\textstyle\bigwedge^rM_0-I)
 \longrightarrow\textstyle\bigwedge^r\Lambda
 \xrightarrow{c_*}H_r(W;\mathbb Z)\longrightarrow0,\qquad
 \operatorname{sp}^r=i_F^*r^*=c^*.
\]
The finite-band attachment has the oriented map
\[
 \alpha_r=(\pi_{1*},-\pi_{2*}):
 H_r(F)\longrightarrow H_r(S_1)\oplus H_r(S_2),
\]
and, after splitting primitive cusp kernels, the only global blocks are
\[
 D_1=\begin{pmatrix}1&0&0\\0&1&0\\-12&0&-p\end{pmatrix},\qquad
 D_2^{\mathrm{red}}=\begin{pmatrix}-p&0\\0&1\end{pmatrix},\qquad
 \det D_1=\det D_2^{\mathrm{red}}=-p.
\]
They are unimodular for $p=-1$ and reproduce the same integral homology.
Dualizing these oriented collar maps gives the three displayed
transgressions, so the Leray and cellular proofs agree map-by-map without
depending on one another.

The explicit candidate and terminal algebra remain attributed to the
Claude/Fabel source.  The audits isolated the missing geometric
identifications, and the present AI-produced verification aid reconstructs
the gluing, both topology certificates, meromorphic descent, Picard theory,
conductor mechanism, Hodge and Bott--Chern calculations, and the CDP repair
as checkable proofs.  Its purpose is to reduce the reconstruction burden for
human experts, not to claim priority for the candidate or to substitute for
independent specialist adjudication.  Root \texttt{main.tex} is the annotated
edition and \texttt{workbench.tex} contains the complete proofs and sources.

\clearpage
\overviewsection{3. Algebraic reduction, Picard, Hodge, and Bott--Chern invariants}

For a very general fibre,
\[
 \operatorname{NS}(F_b)=\mathbb Z\eta,\qquad
 \eta=u\wedge w+6\gamma\wedge\delta,\qquad
 \eta^2=12\gamma\wedge u\wedge w\wedge\delta.
\]
Because $\det_{\mathbb R}\Pi<0$, complex and marked orientations are opposite:
\[
 \int_{(F_b,\mathrm{or}_{\rm mark})}\eta^2=12,\qquad
 \int_{(F_b,\mathrm{or}_{\mathbb C})}\eta^2=-12.
\]
The Appell--Humbert form has signature $(1,1)$, hence $a(F_b)=0$; the
meromorphic descent theorem gives
\[
 \mathcal M(X)=f^*\mathbb C(t),\qquad a(X)=1.
\]
Writing $\mathfrak K=f_*\omega_{X/B}$, one has
\[
 \mathfrak K\simeq\mathcal O_B(1),\quad
 f_*\mathcal O_X=\mathcal O,\quad
 R^1f_*\mathcal O_X=\mathcal O\oplus\mathcal O(-1),\quad
 R^2f_*\mathcal O_X=\mathfrak K^\vee=\mathcal O(-1),
\]
and $\bigwedge^2R^1f_*\mathcal O_X\simeq R^2f_*\mathcal O_X$.
The full canonical-bundle chain is
\[
 \begin{gathered}
 L:=K_X\otimes\mathcal O_X(-2S_2),\qquad
 M:=f_*L,\qquad f^*M\simeq L,\qquad
 \mathfrak K=M(2),\qquad M\simeq\mathcal O_B(-1),\\
 K_X\simeq f^*\mathcal O_B(-1)\otimes\mathcal O_X(2S_2),\qquad
 K_X^{\otimes4}\simeq f^*\mathcal O_B(-2).
 \end{gathered}
\]
Leray and the exponential sequence give
\[
 H^1(X,\mathcal O_X)=\mathbb C,\quad H^2(X,\mathcal O_X)=0,\quad
 \operatorname{Pic}(X)\simeq H^1(X,\mathcal O_X)\simeq\mathbb C.
\]

For $D=W+S_1+S_2$, $\Sigma=\{p_0,p_1,p_2\}$ and
$\mathcal E=\Omega_X^1(\log D)/f^*\Omega_B^1(\log\Sigma)$, the cyclic
invariant lattices give
\[
 f_*\mathcal E=\mathcal O\oplus\mathcal O(-1),\quad
 R^1f_*\mathcal E=\mathcal O^{\oplus2}\oplus\mathcal O(-1)^{\oplus2},\quad
 R^2f_*\mathcal E=\mathcal O\oplus\mathcal O(-1).
\]
The typed Kodaira--Spencer maps
\[
 \kappa_0:f_*\mathcal E\to
 \Omega_B^1(\log\Sigma)\otimes R^1f_*\mathcal O_X,\qquad
 \kappa_1:R^1f_*\mathcal E\to\mathcal O_B
\]
are certified by the load-bearing local equations
\[
 \begin{gathered}
 \operatorname{div}(\det\kappa_0)=p_1+p_2,\qquad
 \operatorname{coker}\kappa_0=\mathbb C_{p_1}\oplus\mathbb C_{p_2},\qquad
 \phi\longmapsto(\phi(p_0),3\phi(p_1),4\phi(p_2)),\\
 \rho_*=(0,1,-1,0):\mathbb C^4\to\mathbb C\theta,\qquad
 \kappa_1(h\otimes\gamma)(p_2)=-\tfrac14\dot\mu(z_2)\ne0.
 \end{gathered}
\]
The first row is the normalized residue cancellation, the second the cusp
Čech sign and invariant order-four scalar.  Thus
\[
 f_*\Omega_X^1=\mathcal O(-2),\qquad
 R^1f_*\Omega_X^1=\mathcal O\oplus\mathcal O(-1)^{\oplus2},\qquad
 R^2f_*\Omega_X^1=\mathcal O\oplus\mathcal O(-1),
\]
and coherent Leray yields
\[
\begin{array}{c|cccc}
q\backslash p&0&1&2&3\\ \hline
3&0&0&0&1\\2&0&1&2&1\\1&1&2&1&0\\0&1&0&0&0.
\end{array}
\]
The two nontrivial Fr\"olicher rows are
$0\to\mathbb C\to\mathbb C^2\to\mathbb C\to0$;
$\dim E_1=10$, $\dim E_2=2$, and $E_2=E_\infty$.  McHugh's formulas give
\[
 (h^k_{BC})=(1,0,2,2,2,2,1),\quad
 (h^k_A)=(1,2,2,2,2,0,1),\quad
 (\Delta^k)=(0,2,4,4,4,2,0).
\]

Let $\mathfrak g=[G]_{\bar\partial}$ be induced by $\gamma$, and choose a
coherent-Leray lift $\mathfrak q=[Q]_{\bar\partial}$ of $\eta$.  In complex
orientation the decisive fibre pairings are
\[
 \int_{F_b}Q_F^2=-12,\qquad
 \int_{F_b}Q_F\wedge\Gamma\wedge d\zeta_2=-1.
\]
For $\exp:H^1(X,\mathcal O_X)\to\operatorname{Pic}(X)$ induced by
$z\mapsto e^{2\pi iz}$,
\[
 H^0(B,R^1f_*\mathbb Z)=\mathbb Z(12\gamma),\qquad
 \exp(12\mathfrak g)=f^*\mathcal O_B(1).
\]
Put $\mathsf A=[\partial G]_{BC}$,
$\mathsf B=[\bar\partial\bar G]_{BC}$.  In a unit frame,
\[
 \bar\partial_zs=-2\pi izG\otimes s,\qquad
 \nabla_z=d-2\pi i(zG+\bar z\,\bar G),\qquad
 c_{1,BC}(\exp(z\mathfrak g))=z\mathsf A+\bar z\,\mathsf B,
\]
so
$c_{1,BC}(f^*\mathcal O_B(1))=12(\mathsf A+\mathsf B)$.
If $N$ is spanned by
\[
 \mathsf A,\mathsf B,\ [\partial Q]_{BC},[\bar\partial\bar Q]_{BC},\
 [\partial(G\wedge Q)]_{BC},[\bar\partial(\bar G\wedge\bar Q)]_{BC},\
 [\partial(Q^2)]_{BC},[\bar\partial(\bar Q^2)]_{BC},\
 [\nu_{\rm top}]_{BC},
\]
with $\int_X\nu_{\rm top}=1$, then
\[
 \mathbf1=[1]_{BC},\qquad
 H^{*,*}_{BC}(X)=\mathbb C\mathbf1\oplus N,\qquad N^2=0.
\]
After fixing $\mathfrak g$ and one coherent-Leray lift $\mathfrak q$, smooth
gauge changes add Dolbeault coboundaries and preserve all displayed classes,
pairings, nonvanishing arguments and products.  The invariant-basis,
residue, Čech and scalar computations are in
\path{supporting_materials/workbench/research/hodge_frolicher.tex};
the product proofs are in
\path{supporting_materials/workbench/research/bott_chern_algebra.tex}.

\overviewsection{4. CDP20, the analytic Picard certificate, and the corrected theorem}

For a reduced nonnormal fibre $S$, put
$\mathcal K_S=\ker(\Omega_X^1|_S\to\widetilde\Omega_S^1)$.  The exact
sequence and the group surviving the two vanishings used in CDP20 are
\[
 0\to N^*_{S/X}\to\mathcal K_S\to\operatorname{tors}\Omega_S^1\to0,
\]
\[
 H^0(S,\Omega_X^1|_S\otimes A)=
 \ker\!\left\{
 H^0(S,\operatorname{tors}\Omega_S^1\otimes A)
 \xrightarrow{\delta_A}H^1(S,N^*_{S/X}\otimes A)\right\}.
\]
CDP20 does not prove $\delta_A$ injective.  Its second imported step must
also use coinvariants:
\[
 (K\rtimes Q)^{\rm ab}\simeq
 Q^{\rm ab}\oplus(K^{\rm ab})_Q.
\]
The contradiction to CDP20 Theorem~2.2 and its $S^6$ corollary
recorded in the workbench uses its complete analytic construction,
topology and algebraic-reduction proofs, giving $a(X)=1$ and
$c_3(X)=2$; it is not inferred from the two proof-route defects above.

It remains to prove the Picard restriction needed by the cusp conductor
without recycling global topology.  For
\[
 E=c_1(L|_F)=kq\in\operatorname{Alt}^2(\Lambda,\mathbb Z),\qquad
 V_{\mathbb R}=\operatorname{Hom}(\Lambda,\mathbb R),
\]
the global factors give
\[
 d:\pi_1(P)\to V_{\mathbb R},\qquad
 d(gh)=d(g)+T_gd(h),\qquad
 \bar d_1+\bar d_2+\bar d_0=0\quad\text{in }V_G\otimes\mathbb R.
\]
Distinguish every finite covering level:
\[
 \widetilde J_j=\Delta_j\times\mathbb C^2
 \xrightarrow{\,r_j\,}J'_j,\qquad
 q_{j,\bullet}\circ r_j=r_j\circ Q_{j,\bullet},
 \quad \bullet\in\{\log,\mathrm{lin}\},
\]
\[
 \begin{aligned}
 Q_{j,\log}\tau_\lambda Q_{j,\log}^{-1}&=\tau_{A_j\lambda},&
 (Q_{j,\log})^{m_j}&=\tau_{v_j},\\
 Q_{j,\mathrm{lin}}\tau_\lambda Q_{j,\mathrm{lin}}^{-1}
 &=\tau_{A_j\lambda},&
 (Q_{j,\mathrm{lin}})^{m_j}&=1.
 \end{aligned}
\]
On $J_j^{\prime *}=J'_j|_{\Delta_j^*}$,
\[
 \begin{gathered}
 h_jq_{j,\log}h_j^{-1}=q_{j,\mathrm{lin}},\\
 p_{j,\log}:J_j^{\prime *}\to
 N_j^*:=J_j^{\prime *}/\langle q_{j,\log}\rangle,\qquad
 p_{j,\mathrm{lin}}:J_j^{\prime *}\to
 J_j^*:=J_j^{\prime *}/\langle q_{j,\mathrm{lin}}\rangle,\\
 \bar h_j:N_j^*\xrightarrow{\sim}J_j^*,\qquad
 p_{j,\mathrm{lin}}h_j=\bar h_jp_{j,\log},\qquad
 p_{j,\log}^*(L|_{N_j^*})
 \simeq h_j^*p_{j,\mathrm{lin}}^*(L|_{J_j^*}).
 \end{gathered}
\]
Put
\[
 \widetilde J_j^*=\Delta_j^*\times\mathbb C^2,\quad
 \Phi_{j,\log}=p_{j,\log}r_j,\quad
 \Phi_{j,\mathrm{glob}}=p_{j,\mathrm{lin}}h_jr_j
 =\bar h_j\Phi_{j,\log}.
\]
The pulled gluing isomorphism
\[
 \Theta_j:\Phi_{j,\log}^*(L|_{N_j^*})
 \xrightarrow{\sim}\Phi_{j,\mathrm{glob}}^*(L|_{J_j^*})
\]
intertwines descent and induces a centre-preserving isomorphism of the two
$\Gamma_j^{\log}$-factor extensions, where
\[
 \Gamma_j^{\log}=
 \langle\tau_\lambda,Q_{j,\log}:\lambda\in\Lambda\rangle.
\]
For normalized vertical lifts
\[
 U_\lambda^{(j)}U_\mu^{(j)}
 =\mathbf c\!\left(\tfrac12E(\lambda,\mu)\right)U_{\lambda+\mu}^{(j)},
\]
define
\[
 \widetilde G_\alpha U_{A_\alpha^{-1}\lambda}^{(j)}
 \widetilde G_\alpha^{-1}
 =\mathbf c(\langle\widetilde d_j(\alpha),\lambda\rangle)
 U_\lambda^{(j)}.
 \tag{5.17-A14c}
\]
Then
\[
 \widetilde d_j(\alpha\beta)=
 \widetilde d_j(\alpha)+T_\alpha\widetilde d_j(\beta).
\]
Compatible basepoints and
$\ell'(w)=\ell(g_j^{-1}w)-2\pi i/m_j$ identify the clockwise zero-section
endpoint with exactly $Q_{j,\log}$, not an inverse or lattice translate.
Thus, for a comparison character $a_j$,
\[
 \widetilde d_j(Q_{j,\log})=d_j+(T_j-I)a_j,\qquad
 \widetilde d_j(\tau_v)=\iota_vE.
\]
Since $N_{T_j}(T_j-I)=0$,
\[
 N_{T_j}d_j
 =N_{T_j}\widetilde d_j(Q_{j,\log})
 =\widetilde d_j((Q_{j,\log})^{m_j})
 =\widetilde d_j(\tau_{v_j})
 =\iota_{v_j}E.
 \tag{5.17-A15}
\]
The comparison is independent of frames, lifts, normalized vertical
systems, and logarithm branch; branch change
$H_j\mapsto\tau_{nv_j}H_j$ changes $a_j$ by
$n\iota_{v_j}E$, which is $T_j$-invariant.

At the cusp the argument principle gives $\bar d_0=0$.  With
\[
 \eta_1=2u+w+3\delta,\qquad \eta_2=u+w+2\delta,
\]
\[
 \iota_{v_1}q=2\eta_1,\quad
 \iota_{v_2}q=-3\eta_2,\quad
 [\eta_1]=3[\delta],\quad[\eta_2]=2[\delta],
\]
equation (5.17-A15) yields
\[
 \bar d_1=2k[\delta],\qquad
 \bar d_2=-\tfrac32k[\delta],\qquad
 0=\bar d_1+\bar d_2+\bar d_0=\tfrac{k}{2}[\delta].
\]
Hence $k=0$ and $E=0$.

For central triviality, put
\[
 M=L|_{N_0},\qquad \mathcal M=\pi^*M,\qquad
 \mathcal R=\mathcal M\otimes\mathcal O(-D_a).
\]
The negative-divisor linearization is the inverse-dual
\[
 \mu_\lambda^-=((\mu_\lambda^+)^{-1})^\vee,\qquad
 \theta_\lambda^{\mathcal R}
 =\theta_\lambda^{\mathcal M}\otimes\mu_\lambda^-.
\]
In the stated lattice bases,
\[
 B_0:\bar\Lambda\xrightarrow{\sim}\Lambda_{\rm tor},\qquad
 [B_0]=\begin{pmatrix}0&1\\-1&0\end{pmatrix}.
\]
If $A=\nu^*(M|_W)$ has degrees
$(d_1,d_2,d_3,d_1,d_2,d_3)$ and support matrix
\[
 [B]=\begin{pmatrix}d_2+d_3&d_3\\d_3&d_1+d_3\end{pmatrix},
\]
then
\[
 E_M((\lambda,0),(0,m))=-B(B_0\lambda,m),\qquad
 E_M=c_1(M|_{F_t})=c_1(L|_F)=E=0.
\]
Unimodularity of $B_0$ gives $B=0$, hence
$d_1=d_2=d_3=0$ and $A\simeq\mathcal O_{\mathrm{dP}_6}$.  Therefore
$\nu^*i^*:\operatorname{Pic}(X)\to\operatorname{Pic}(\mathrm{dP}_6)$ is
zero.

For $F_L=T_X\otimes L$ and
$A_L=(L^\vee\otimes\omega_X)|_W$, duality and base change give
\[
 H^2(W,F_L|_W)^\vee\simeq
 H^0(W,\Omega_X^1|_W\otimes A_L),\qquad
 (R^2f_*F_L)_{p_0}\otimes\mathbb C\simeq H^2(W,F_L|_W).
\]
The conductor section and the just-proved Picard restriction show
$R^2f_*(T_X\otimes L)\ne0$ for every $L\in\operatorname{Pic}(X)$.

A corrected CDP statement is: if one $L_0$ satisfies
$H^2(X,T_X\otimes L_0)=0$, then for general $M\in\operatorname{Pic}(X)$,
\[
 H^i(X,T_X\otimes M)=0\ (i\ne1),\qquad
 \tfrac12\langle c_3(T_X),[X]\rangle
 =-h^1(X,T_X\otimes M)\le0.
\]
Relatively, one may instead assume one
$R^2f_*(T_X\otimes L)=0$; for reduced full fibres, the exact replacement
for the invalid local step is injectivity of the conductor connecting map
$\delta_c$ on every whole fibre.  The candidate fails this relative
sufficient condition by the preceding conductor nonvanishing.  Separately,
because $c_3(X)=2$, the contrapositive of the corrected global statement gives
$H^2(X,T_X\otimes L_0)\ne0$ for every $L_0$.  Full local calculations are in
\path{supporting_materials/workbench/research/analytic_picard_completion.tex}
and
\path{supporting_materials/workbench/research/cdp_program.tex}.

\clearpage
\overviewsection{5. New cusp mathematics: Kummer extension, sphere carriers, and finite gauge systems}

The subsequent work adds several precise results which are logically separate
from the construction and recognition of $X$.  They concern the local toric
cusp, its metric limits, topological carriers for its expanding quotients,
and a finite classical/Gaussian gauge model.  None changes the provenance
statement in Section~2: these are reconstruction and extension results in the
AI-produced workbench, not claims of priority for the original candidate.

\medskip\noindent
\textit{The local toric Kummer map and its two degrees.}
Write the completed cusp as
$f_0:N_0=Y/\bar\Lambda\to\Delta_t$, with dense-torus
action
\[
 \Psi_{\bar\lambda}(x,t)=
 \bigl(e^{2\pi iC(t)\bar\lambda}t^{B_0\bar\lambda}x,t\bigr),
 \qquad B_0=\begin{pmatrix}0&1\\-1&0\end{pmatrix}.
\]
For an integer $N\ge1$, let
$\rho_N:\Delta_v\to\Delta_t$ be $\rho_N(v)=v^N$, take a copy
$Y^{(N)}$ of the same toric variety, and replace the source multiplier by
the holomorphic root
\[
 C_N(v)=\frac1N C(v^N),\qquad
 \Psi^{(N)}_{\bar\lambda}(\xi,v)=
 \bigl(e^{2\pi iC_N(v)\bar\lambda}
 v^{B_0\bar\lambda}\xi,v\bigr).
\]
Its quotient and reduced central fibre are
\[
 f_N:N_0^{\langle N\rangle}:=Y^{(N)}/\bar\Lambda\longrightarrow\Delta_v,
 \qquad
 W_N:=\bigl(f_N^{-1}(0)\bigr)_{\rm red}.
\]
Then the toric power map
\[
 P_N(\xi_1,\xi_2,v)=(\xi_1^N,\xi_2^N,v^N)
\]
is equivariant and descends to a finite flat map
\[
 \kappa_N:N_0^{\langle N\rangle}\longrightarrow N_0,
 \qquad f_0\circ\kappa_N=\rho_N\circ f_N,
 \qquad \deg\kappa_N=N^3,\qquad
\kappa_N^*W=NW_N,\qquad
 \operatorname{Ram}(\kappa_N)=(N-1)W_N.
\]
The exponent $3$ is the degree over the \emph{unchanged} cusp disc: two
fibre coordinates and the base coordinate are all raised to the $N$th
power.  After replacing the target by the base change
$N_0[N]=N_0\times_{\Delta_t}\Delta_v$, the factorization
\[
 N_0^{\langle N\rangle}\xrightarrow{\,Q_N\,}N_0[N]\longrightarrow N_0
\]
has generic degree $N^2$.  For $N>1$, $Q_N$ is quasi-\'etale but not flat;
its non-\'etale and nonflat loci lie exactly over the double and triple strata.
At a triple point its closed-fibre length is
$3N^2-3N+1>N^2$, and at a double stratum the transverse closed-fibre length
is $2N-1>N$.  On the reduced central fibre, the generic degree is $N^2$ and
after normalizing the reduced source and target central fibres, the induced
map is toric multiplication
$[N]:\mathrm{dP}_6\to\mathrm{dP}_6$, ramified with index $N$ along all six
boundary curves.  For $v\ne0$, the relative fibre map $Q_{N,v}$ has deck
group $(\mathbb Z/N\mathbb Z)^2$, and
\[
S_N^{-1}M_0^NS_N=M_0,\qquad S_N=\operatorname{diag}(1,1,N,N).
\]
This is a completed \emph{local} cusp construction.  It is not a
degree-$N^2$ self-map over the original disc, not a normalization of the
base change, and not yet a cover over the full three-point base.

\medskip\noindent
\textit{Native and auxiliary pointed limits.}
Put $t_c=e^{-2\pi T}$ on the positive radial ray and write the marked period
matrix and base point as
\[
 P_T=(v_{\gamma,T}\ v_{u,T}\ e_1\ e_3),\qquad
 F_T=\mathbb R^4/P_T\mathbb Z^4,\qquad o_T=[0].
\]
Two primitive periods remain of length one, while the other two escape
without hidden short integer cancellations: for constants
$c_-,c_+>0$ independent of $T$,
\[
 \rho_T:=\min_{(a,b)\ne(0,0)}
 \operatorname{dist}\!\left(
 av_{\gamma,T}+bv_{u,T},
 \mathbb Z\langle e_1,e_3\rangle\right)
 \ge T-c_-.
\]
Consequently the exact native limit of the unmodified four-torus fibres
\(F_T\), in the unscaled Euclidean quotient metric, is
\[
 (F_T,o_T)\longrightarrow
 (\mathbb T^2\times\mathbb R^2,(0,0)),
\]
and every fixed-radius ball is eventually isometric to the corresponding
ball in that product.  An $\mathbb R^4$ pointed limit of those four-torus
fibres themselves appears only after the added degree-$M^2$ local cover
obtained by replacing the two short periods by $M$ times those periods:
\[
 \Lambda_{T,M}=
 \mathbb Z\langle v_{\gamma,T},v_{u,T},Me_1,Me_3\rangle,
 \qquad
 \widetilde F_{T,M}=\mathbb R^4/\Lambda_{T,M},
\]
\[
 \min\{T-c_-,M\}\leq\operatorname{sys}(\Lambda_{T,M})
 \leq\min\{T+c_+,M\}.
\]
For every integer-valued $M(T)\in\mathbb Z_{>0}$ with
$M(T)\to\infty$, the pointed cover tori converge locally isometrically and
in pointed Gromov--Hausdorff topology to $\mathbb R^4$; on each eventual
Euclidean base-point neighbourhood this convergence is smooth.  This
auxiliary cover is the regular-fibre restriction of the relative
degree-$M^2$ Kummer factor above; the completed map to the original cusp
family still has degree $M^3$.

\medskip\noindent
\textit{Exact scalar counting and a fibrewise magnetic gap.}
Let
\[
 D_T:=-\det_{\mathbb R}P_T
 =\operatorname{Vol}(F_T)=T^2+O(T),
 \qquad
 N_T(E):=\#\{k\in\mathbb Z^4:
 4\pi^2\lVert P_T^{-\mathsf T}k\rVert^2\le E\}.
\]
Counting Fourier modes with multiplicity gives, for every fixed $E\ge0$,
\[
 N_T(E)=\pi D_T\sum_{n\in\mathbb Z^2}
 \left(\frac{E}{4\pi^2}-\lVert n\rVert^2\right)_++O_E(T),
 \qquad (x)_+:=\max\{x,0\}.
\]
In particular, below the first transverse threshold $0<E<4\pi^2$,
\[
 N_T(E)=\frac{ED_T}{4\pi}+O_E(T)
 =\frac{ET^2}{4\pi}+O_E(T).
\]
Thus the fixed-energy population has a two-dimensional leading density,
even though every individual compact fibre has discrete spectrum.

There is also an exact full-four-torus statement which does not pass through
an elliptic quotient.  Fix a nondegenerate integral alternating form
$\mathsf E=(E_{jk})$ on the marked lattice, let
\[
 d:=|\operatorname{Pf}(\mathsf E)|
 =|E_{12}E_{34}-E_{13}E_{24}+E_{14}E_{23}|>0,
\]
and let $\Delta_{\mathsf E,T}$ be the Friedrichs magnetic Laplacian for the
constant-curvature connection with curvature $2\pi i\mathsf E$.  If
\[
 \operatorname{Spec}\!\left(
 (P_T^{\mathsf T}P_T)^{-1/2}(2\pi\mathsf E)
 (P_T^{\mathsf T}P_T)^{-1/2}\right)
 =\{\pm ib_1,\pm ib_2\},\qquad b_1,b_2>0,
\]
then
\[
 \operatorname{Spec}(\Delta_{\mathsf E,T})
 =\{(2n_1+1)b_1+(2n_2+1)b_2:n_1,n_2\in\mathbb Z_{\ge0}\},
\]
\[
 b_1b_2=\frac{4\pi^2d}{D_T},\qquad
 \operatorname{gap}(\Delta_{\mathsf E,T})
 =2\min\{b_1,b_2\}
 \le4\pi\sqrt{\frac d{D_T}}=O(T^{-1}).
\]
This magnetic theorem is fibrewise for a fixed marked integral flux; it
asserts neither full-monodromy invariance of that flux nor a Yang--Mills
mass-gap result.

\medskip\noindent
\textit{Exact $S^{12}$ and $S^{24}$ source maps and their metrics.}
Let
\[
 K_\Lambda=\mathbb Z\langle\widehat w,\widehat\delta\rangle,
 \qquad \overline{\mathbb T}^{\,2}=
 \frac{(\Lambda/K_\Lambda)\otimes\mathbb R}{\Lambda/K_\Lambda}.
\]
Put
\[
 Y_T=
 \begin{pmatrix}
  6\operatorname{Im}\mu_T&\operatorname{Im}\tau_T\\
  \operatorname{Im}\beta_T&\operatorname{Im}\mu_T
 \end{pmatrix},\qquad
 Q_T=\mathbb R^2/Y_T\mathbb Z^2,\qquad
 \iota_T:\overline{\mathbb T}^{\,2}\xrightarrow{\sim}Q_T,
 \quad[y]\longmapsto[Y_Ty].
\]
The quotient $\pi_T:F_T\twoheadrightarrow Q_T$ is the radial restriction
of a quotient over the whole punctured cusp.  Collapsing the central fibre
and the exterior boundary gives the continuous surjective based quotient
denoted $q_2$ in the proof; in this overview write it as
\[
 \rho_{2,\mathrm{cone}}:(X,x_0)\longrightarrow
 (\mathcal A_{\mathrm{cone}},v_0),\qquad
 \mathcal A_{\mathrm{cone}}
 =\operatorname{Susp}(S^1\times\overline{\mathbb T}^{\,2}).
\]
The full-monodromy-invariant character lattice is
$\mathbb Z\varepsilon_\gamma$; fix the displayed primitive generator.  It
gives a circle quotient on the regular base and explicit Seifert targets at
the two finite fillings.  Gluing them gives a compact orientable
three-manifold $\mathcal M_{\mathrm{Seif}}$ with boundary $\mathbb T^2$;
coning that entire boundary gives the map denoted $q_1$ in the proof,
written here as
\[
 \mathcal C_{\mathrm{cone}}=
 \mathcal M_{\mathrm{Seif}}\cup_{\mathbb T^2}
 \operatorname{Cone}(\mathbb T^2),\qquad
 \rho_{1,\mathrm{cone}}:(X,x_0)\longrightarrow
 (\mathcal C_{\mathrm{cone}},c_0).
\]
There is also the cusp-supported target
$\mathcal B_{\mathrm{cone}}=\operatorname{Susp}(S^1\times S^1)$ and a
based quotient $\rho_{1,\mathrm{cusp}}:X\to\mathcal B_{\mathrm{cone}}$.
After a based choice $X\cong_{\rm diff}S^6$, basedness induces
\[
 \rho_{2,\mathrm{cone}}^{\wedge2}:
 S^{12}\cong X^{\wedge2}\longrightarrow
 \mathcal A_{\mathrm{cone}}^{\wedge2},
 \qquad
 \rho_{1,\mathrm{cone}}^{\wedge4}:
 S^{24}\cong X^{\wedge4}\longrightarrow
 \mathcal C_{\mathrm{cone}}^{\wedge4}.
\]
The corresponding Cartesian restrictions, before passage to the smash
quotient and after identifying their radial images, are exactly
\[
 \left.\rho_{2,\mathrm{cone}}^{\times2}\right|_{F_T^{\times2}}
 =\pi_T^{\times2}:F_T^{\times2}\to Q_T^{\times2},
 \qquad
 \left.\rho_{1,\mathrm{cone}}^{\times4}\right|_{F_T^{\times4}}
 =\chi_{\gamma,T}^{\times4}:F_T^{\times4}\to C_T^{\times4},
\]
where
\[
 C_T=\mathbb R/\ell_T\mathbb Z,\qquad
 \ell_T=T+H_{\mathrm{cusp}}-B_{\mathrm{cusp}}+O(T^{-1}),\qquad
 H_{\mathrm{cusp}}=\operatorname{Im}h(0),\quad
 B_{\mathrm{cusp}}=\operatorname{Im}b(0).
\]

Here are the target-slice embeddings, with no implicit identification.  Put
$r_T=e^{-2\pi T}$ and
\[
 \vartheta_T:\mathbb R/\mathbb Z\xrightarrow{\sim}C_T,
 \qquad[a]\longmapsto[\ell_Ta].
\]
In the suspension/cone coordinates used above,
\[
 \begin{aligned}
 j_T(y_1,y_2)
 &=\bigl[[r_T,1,\iota_T^{-1}(y_1)],
          [r_T,1,\iota_T^{-1}(y_2)]\bigr]_{\wedge}
 \in\mathcal A_{\mathrm{cone}}^{\wedge2},\\
 k_T^{\mathcal B}(c_1,\ldots,c_4)
 &=\bigl[[r_T,1,\vartheta_T^{-1}(c_1)],\ldots,
          [r_T,1,\vartheta_T^{-1}(c_4)]\bigr]_{\wedge}
 \in\mathcal B_{\mathrm{cone}}^{\wedge4},\\
 k_T^{\mathcal C}(c_1,\ldots,c_4)
 &=\bigwedge_{j=1}^4
   [r_T,(1,\vartheta_T^{-1}(c_j))]
 \in\mathcal C_{\mathrm{cone}}^{\wedge4}.
 \end{aligned}
\]
Each is a topological embedding.  With the intrinsic product quotient
metrics,
\[
 (Q_T^{\times2},0)\longrightarrow(\mathbb R^4,0),
 \qquad
 (C_T^{\times4},0)\longrightarrow(\mathbb R^4,0)
\]
pointed locally isometrically and in pointed Gromov--Hausdorff topology,
with smooth convergence on the eventual Euclidean base-point
neighbourhoods.

These copy numbers are sharp for the unrescaled quotient norm induced by
the marked Euclidean fibre metric.  For a connected torus quotient of $r$
copies over the punctured cusp, divergence of the quotient systole forces
the bounded lattice $K_\Lambda^{\oplus r}$ into the kernel, so the quotient
rank is at most $2r$; rank four requires $r\ge2$ and is attained by
$Q_T^{\times2}$.  For a quotient over the full regular base, even after a
finite base cover, monodromy also forces
$\widehat u^{\oplus r}$ into the primitive kernel, so the quotient rank is
at most $r$; rank four requires $r\ge4$ and is attained by
$C_T^{\times4}$.  Hence $S^{12}$ and $S^{24}$ are the smallest smash-sphere
sources supplied by the native cusp-local and full-base quotient routes,
respectively.

There are also genuine sphere targets.  Collapsing the one-skeleton
\[
 K^{(1)}=((\mathbb R/\mathbb Z)\times\{0\})
 \cup(\{0\}\times(\mathbb R/\mathbb Z))
 \subset\overline{\mathbb T}^{\,2}
\]
and forgetting the base angle give continuous surjective based quotient
maps $\widehat\rho_3:X\to S^3$ and
$\widehat\rho_\gamma:X\to S^2$.  Their smash powers are
\[
 \widehat\rho_3^{\wedge2}:S^{12}\twoheadrightarrow S^6,
 \qquad
 \widehat\rho_\gamma^{\wedge4}:S^{24}\twoheadrightarrow S^8.
\]
Here \(\widehat\rho_3=\widehat q_3\) and
\(\widehat\rho_\gamma=\widehat q_\gamma\) in the detailed construction.
Both full sphere carriers are based null-homotopic, by two distinct exact
calculations.  For every based map \(f:S^6\to S^2\),
\[
 S^{12}=S^6\wedge S^6
 \xrightarrow{\ f\wedge\operatorname{id}_{S^6}=\Sigma^6f\ }
 S^2\wedge S^6=S^8
 \xrightarrow{\ \operatorname{id}_{S^2}\wedge f\ }
 S^2\wedge S^2=S^4.
\]
Since
\(\pi_{12}(S^8)\cong\pi_4^{\rm S}=0\), the first arrow is based
null-homotopic, and hence so is \(f\wedge f\).  Therefore
\[
 [f^{\wedge r}]=0\in\pi_{6r}(S^{2r})
 \qquad\text{for every }r\ge2,
\]
because
\(f^{\wedge r}=(f\wedge f)\wedge f^{\wedge(r-2)}\).  Taking
\(f=\widehat\rho_\gamma\) proves that the entire surjective quotient
\(\widehat\rho_\gamma^{\wedge4}:S^{24}\twoheadrightarrow S^8\) is based
null-homotopic.

For every based map \(g:S^6\to S^3\), the exact smash factorization is
\[
 g\wedge g=(\Sigma^3g)\circ(\Sigma^6g):
 S^{12}\longrightarrow S^9\longrightarrow S^6.
\]
Let \(\operatorname{st}\) denote stabilization.  The exact low-stem data are
\[
 \pi_6(S^3)\cong\mathbb Z/12,
 \quad
 \pi_3^{\rm S}=\mathbb Z/24\langle\nu\rangle,
 \quad
 \pi_6^{\rm S}=\mathbb Z/2\langle\nu^2\rangle,
\]
\[
 \pi_{12}(S^6)=\mathbb Z/2\langle\nu_6\circ\nu_9\rangle,
 \qquad
 \operatorname{st}:\pi_{12}(S^6)\xrightarrow{\ \cong\ }\pi_6^{\rm S}.
\]
Here the full group table is Hatcher, p.~339, and the signed smash formula
and unstable stabilization calculation are Toda, Propositions~3.1 and~5.11.
Since \(12[g]=0\), one has \(\operatorname{st}([g])=2a\nu\) for some
\(a\in\mathbb Z/12\).  Multiplicativity gives
\[
 \operatorname{st}([g\wedge g])
 =\operatorname{st}([g])^2
 =4a^2\nu^2=0\in\pi_6^{\rm S}.
\]
The stabilization isomorphism therefore gives
\([g\wedge g]=0\in\pi_{12}(S^6)\).
Thus \(g^{\wedge r}\simeq *\) for every \(r\ge2\), and taking
\(g=\widehat\rho_3\) proves
\[
 [\widehat\rho_3^{\wedge2}]=0\in\pi_{12}(S^6),
 \qquad
 [\widehat\rho_\gamma^{\wedge4}]=0\in\pi_{24}(S^8).
\]
The one-copy map \(\widehat\rho_3\) is now proved null-homotopic by
the explicit framed filling described above. The one-copy class of
\(\widehat\rho_\gamma\) remains open. The unstable classes of the
conical-target maps
\(\rho_{2,\mathrm{cone}}^{\wedge2}\),
\(\rho_{1,\mathrm{cusp}}^{\wedge4}\), and
\(\rho_{1,\mathrm{cone}}^{\wedge4}\) also remain open: nullity after
postcomposition with the angle-forgetting quotient does not imply nullity of
the maps before that postcomposition.
For each sufficiently large $T$, their targets contain moving
topologically embedded four-spheres, equipped separately with their
intrinsic quotient length metrics,
\[
 \mathbb S_{12,T}^4=S_T^2\wedge S_T^2\hookrightarrow S^6,
 \qquad
 \mathbb S_{24,T}^4=C_T^{\wedge4}\hookrightarrow S^8,
\]
where $S_T^2=\overline{\mathbb T}^{\,2}/K^{(1)}$ carries the quotient
length metric pulled back from $Q_T$.  Pointing these spaces at the
midpoints $\zeta_{12,T},\zeta_{24,T}$, rather than at their smash basepoints,
every
radius-$R$ ball is exactly a Euclidean four-ball once
$T-c_->4R$, respectively once $\ell_T>4R$.  Thus both moving four-spheres
converge to $\mathbb R^4$ pointed locally isometrically and in pointed
Gromov--Hausdorff topology, again smoothly on the eventual Euclidean
neighbourhoods.  The based inclusion maps
\[
 \mathbb S_{12,T}^4\hookrightarrow S^6,
 \qquad
 \mathbb S_{24,T}^4\hookrightarrow S^8
\]
are null-homotopic: sliding the fixed suspension coordinate to an endpoint
contracts every factor through based maps.  Their four-dimensional
fundamental classes therefore die in the ambient targets.  These two
slice-inclusion null-homotopies are separate from the global equalities
\([\widehat\rho_3^{\wedge2}]=0\in\pi_{12}(S^6)\) and
\([\widehat\rho_\gamma^{\wedge4}]=0\in\pi_{24}(S^8)\).

Deleting the common smash basepoint gives literal nonclosed topological
four-cells in the sphere targets:
\[
 (Q_T\setminus K_T^{(1)})^2
 \xrightarrow{\ \cong\ }
 \mathbb S_{12,T}^4\setminus\{*\}
 \cong\mathbb R^4\hookrightarrow S^6,
\]
\[
 (C_T\setminus\{0\})^4
 \xrightarrow{\ \cong\ }
 \mathbb S_{24,T}^4\setminus\{*\}
 \cong\mathbb R^4\hookrightarrow S^8.
\]
They are open only inside their embedded four-spheres, not in the ambient
sphere, and no fixed compatible compact-target metric makes them globally
isometric to Euclidean $\mathbb R^4$.

The metric scope is exact.  For every unrescaled fixed metric compatible
with any compact cone-target topology, the slices $j_T(Q_T^{\times2})$ and
$k_T^{\mathcal B}(C_T^{\times4})$,
and $k_T^{\mathcal C}(C_T^{\times4})$ converge to the cone base point.
Fixed cusp-adapted
length metrics do admit moving-base-point blow-ups: after factors
$T/r_T$ and $\ell_T/r_T$, the displayed slices converge to
$\mathbb R^4$, while the full conical ambient bubbles are respectively
$\mathbb R^8$ and $\mathbb R^{12}$.  Their conic vertex links have
dimensions $7$ and $11$, so the vertex cone dimensions are $8$ and $12$,
not $4$.  By contrast, an independently chosen smooth Riemannian metric on
the source sphere has ordinary tangent space $\mathbb R^{12}$ or
$\mathbb R^{24}$.  The quotient construction supplies no Riemannian metric
on either source sphere.

Thus $\mathbb R^4$ is neither the whole compact source nor the whole compact
target, and it cannot be a continuous quotient target of a compact sphere.
What is proved is a literal nonclosed topological open cell, an intrinsic
pointed limit of moving flat slices (not Sormani--Wenger
intrinsic-flat convergence) and, for the fixed cusp-adapted target metrics,
a four-plane inside a higher-dimensional moving-point bubble.  No auxiliary
finite-cover unspooling is used in these sphere routes, and no compact
complex structure on $S^{12}$ or $S^{24}$ follows.  Cartesian products also
remain distinct from smash spheres:
$X^{\times2}\simeq S^6\times S^6\ne S^{12}$ and
$X^{\times4}\simeq(S^6)^{\times4}\ne S^{24}$.

The compact carrier subspaces swept out by the synchronized slices are
\[
 \mathcal D_{12}\cong\mathcal D_{24}\cong\operatorname{Cone}(\mathbb T^4).
\]
At their vertices the local integral homology has ranks $4,6,4,1$ in
degrees $2,3,4,5$, respectively.  Thus a vertex is not a
homology-manifold point, and hence not locally a topological manifold (with
or without boundary) of any dimension.  The full smash targets retain still
more topology:
\[
 \mathcal A_{\mathrm{cone}}^{\wedge2}\simeq
 \bigvee^9S^4\vee\bigvee^{18}S^5\vee\bigvee^{15}S^6
 \vee\bigvee^6S^7\vee S^8,
\]
\[
 \mathcal C_{\mathrm{cone}}^{\wedge4}\simeq
 S^8\vee\bigvee^4S^9\vee\bigvee^6S^{10}
 \vee\bigvee^4S^{11}\vee S^{12}.
\]
These are homotopy equivalences, not homeomorphisms.

There is nevertheless a space in which the limiting Euclidean fibre is
literal rather than merely asymptotic.  For $\nu\in\{12,24\}$ put
\[
 A_T^{(12)}=\operatorname{diag}(Y_T,Y_T),\qquad
 A_T^{(24)}=\ell_T I_4,
\]
and quotient $[0,s_0]\times\mathbb R^4$ fibrewise by
$A_{1/s}^{(\nu)}\mathbb Z^4$ for $s>0$, while imposing no identifications
at $s=0$.  The resulting $\mathfrak E_\nu$ is a Hausdorff,
second-countable topological five-manifold with boundary, with
\[
 p_{12}^{-1}(s)\cong Q_{1/s}^{\times2},\quad
 p_{24}^{-1}(s)\cong C_{1/s}^{\times4}\quad(s>0),\qquad
 p_\nu^{-1}(0)\cong\mathbb R^4.
\]
There are continuous surjective comparison maps
$\beta_\nu:\mathfrak E_\nu\twoheadrightarrow\mathcal D_\nu$ which identify
every positive fibre with its radial carrier slice and collapse the entire
$\mathbb R^4$ boundary fibre to the cone vertex.  They are homeomorphisms off
that boundary, but are nonproper and nonquotient; in particular they are not
topological blow-downs or resolutions, and
$\mathcal D_\nu\not\cong\mathfrak E_\nu/\mathbb R^4$ via $\beta_\nu$.
Thus the requested Euclidean object is an intrinsic moving-slice limit and an
actual boundary fibre of a noncompact comparison manifold, not the whole
compact $S^{12}$, $S^{24}$, or compact quotient target.

\medskip\noindent
\textit{Repaired cubical $SU(2)$ flux and the Gaussian theorem.}
The finite model keeps independent integers $M_{\rm cov}\ge1$, the covering
index, and $N_{\rm mesh}\ge2$, the number of cubical subdivisions in each
direction.  Put
\[
 S_{M_{\rm cov}}=\operatorname{diag}(1,1,M_{\rm cov},M_{\rm cov}),
 \qquad P_{T,M_{\rm cov}}=P_TS_{M_{\rm cov}},
\]
\[
 V_{N_{\rm mesh}}=(\mathbb Z/N_{\rm mesh}\mathbb Z)^4,
 \qquad
 E_{N_{\rm mesh}}=V_{N_{\rm mesh}}\times\{1,2,3,4\}.
\]
Recall that $D_T=-\det_{\mathbb R}P_T=T^2+O(T)$, so the cover covolume is
$M_{\rm cov}^2D_T$.  With the normalization
\[
 \langle X,Y\rangle=-\tfrac12\operatorname{tr}(XY),\qquad
 H_{\rm link}=\begin{pmatrix}i&0\\0&-i\end{pmatrix},\qquad
 \lVert H_{\rm link}\rVert^2=1,\qquad
 \exp(2\pi H_{\rm link})=I_2,
\]
the periodic inverse-transport links
\[
 \begin{aligned}
 U_1(n)
 &=\exp\!\left(-\frac{2\pi n_2}{N_{\rm mesh}^2}
 H_{\rm link}\right),\\
 U_2(n)
 &=
 \begin{cases}
  I_2,&0\le n_2\le N_{\rm mesh}-2,\\
  \exp(2\pi n_1H_{\rm link}/N_{\rm mesh}),
      &n_2=N_{\rm mesh}-1,
 \end{cases}\\
 U_3(n)&=U_4(n)=I_2,
 \qquad n\in V_{N_{\rm mesh}}.
 \end{aligned}
\]
give every $(1,2)$ plaquette the same value
$\exp(2\pi H_{\rm link}/N_{\rm mesh}^2)$ and all other positively oriented coordinate
plaquettes the identity.  The principal-log curvature is exactly
$2\pi H_{\rm link}\,dy^1\wedge dy^2$, and the metric-weighted logarithmic action is
\[
 S^{\log}_{T,M_{\rm cov},N_{\rm mesh}}
 =\frac{2\pi^2M_{\rm cov}^2}{g_{\rm YM}^2D_T},
\]
independent of the mesh.  The explicitly defined continuum-matched Wilson
sector has value
\[
 \frac{M_{\rm cov}^2N_{\rm mesh}^4}{2g_{\rm YM}^2D_T}
 \left(2-2\cos\frac{2\pi}{N_{\rm mesh}^2}\right)
 \longrightarrow
 \frac{2\pi^2M_{\rm cov}^2}{g_{\rm YM}^2D_T}.
\]

For $k=(k_1,\ldots,k_4)\in V_{N_{\rm mesh}}$, define the componentwise
Fourier symbol
\[
 \mathfrak d_{N_{\rm mesh}}(k)_j
 =N_{\rm mesh}\left(e^{2\pi ik_j/N_{\rm mesh}}-1\right),
 \qquad 1\le j\le4.
\]
The reduced finite Coulomb space
$\ker d_{N_{\rm mesh}}^*\cap
(\ker\Delta_{N_{\rm mesh}}^{(1)})^\perp$ removes the complete harmonic
$k=0$ fibre.  On that space, Fourier diagonalization of
$\mathcal L=g_{\rm YM}^{-2}d_{N_{\rm mesh}}^*d_{N_{\rm mesh}}$ gives
\[
 \lambda_{T,M_{\rm cov},N_{\rm mesh}}(k)
 =\frac1{g_{\rm YM}^2}
 \left\|P_T^{-\mathsf T}S_{M_{\rm cov}}^{-1}
 \mathfrak d_{N_{\rm mesh}}(k)\right\|^2,
 \qquad k\ne0,
\]
with transverse complex multiplicity $9$.  Its Friedrichs Gaussian
oscillator
\[
 \widehat H^{\rm Gauss}=-\tfrac12\Delta
 +\tfrac12\langle a,\mathcal La\rangle
\]
has a unique Gaussian ground state and exact gap
$\sqrt{\lambda_{\min}(\mathcal L)}$.  The explicit sequence
\[
 T_j^{\rm rad}=j^2,\qquad M_{{\rm cov},j}=j,\qquad
 N_{{\rm mesh},j}=j^3
\]
simultaneously gives a pointed local-isometric $\mathbb R^4$ cover limit,
vanishing mesh size, vanishing Gaussian gap, and vanishing logarithmic
classical action.

This last statement is a finite quadratic-oscillator theorem, not a
four-dimensional interacting Yang--Mills mass-gap result.  Its Hilbert space
is $L^2$ of a noncompact linear Coulomb slice with Lebesgue measure, whereas
compact-link lattice theory uses Haar measure on
$SU(2)^{E_{N_{\rm mesh}}}$ followed by the nonlinear gauge constraint.
The exponential is only a local chart, no
unitary identification of these Hilbert spaces is constructed, a link
configuration defines a Dirac distribution rather than an $L^2$ state, and
there is no nonlinear Hamiltonian, transfer matrix, renormalized continuum
limit, or Osterwalder--Schrader reconstruction in the result.

\clearpage
\overviewsection{6. Split \(SU(r)\), the \(SU(3)\) Gram and Lorentz
carriers, and the exact Erd\H{o}s--Straus--Koide incidence bridge}

The 3--4 September 2026 continuation separates the structures that first
appeared together as a heuristic and then computes the strongest exact maps
among them: a gauge field on each native four-torus; the two- and three-column
\(SU(3)\) orbit cones; a split Lorentz observation attached to the toroidal
angular epimorphism; the full retained four-vector--spin-representation--scalar
fibre; and the arithmetic incidence correspondence.  Each now has a complete
theorem, but the maps
which would identify them with one another either fail by an explicit
invariant or have not been constructed.

\medskip\noindent
\textit{The full split-Cartan \(SU(r)\) theorem on the native fibres.}
Fix \(r\ge2\) and a nonzero integral weight
\[
 \mathbf m=(m_1,\ldots,m_r),\qquad
 \sum_i m_i=0,\qquad s_2(\mathbf m)=\sum_i m_i^2.
\]
For the full-flux line bundle \(\mathscr L_{\eta,T}\to F_T\), put
\[
 \mathscr E_{\mathbf m,T}
 =\bigoplus_{i=1}^r\mathscr L_{\eta,T}^{\,m_i},
 \qquad
 H_{\mathbf m}=i\operatorname{diag}(m_1,\ldots,m_r).
\]
The ordered splitting gives a parallel determinant trivialization, and
\[
 c_1(\mathscr E_{\mathbf m,T})=0,\qquad
 c_2(\mathscr E_{\mathbf m,T})
 =-\frac{s_2(\mathbf m)}2[\eta_T]^2.
\]
Since the marked and physical complex orientations are opposite,
\[
 \left\langle c_2,[F_T]_{\rm mark}\right\rangle=-6s_2(\mathbf m),
 \qquad
 k_{\mathbf m}:=
 \left\langle c_2,[F_T]_{\rm phys}\right\rangle=6s_2(\mathbf m).
\]
The congruence \(s_2(\mathbf m)\equiv\sum_i m_i=0\pmod2\) shows that this
split family realizes precisely a subset of the sectors
\(k_{\mathbf m}\in12\mathbb Z_{>0}\).

The diagonal connection and curvature are
\[
 \nabla^{\mathbf m,T}
 =d-2\pi(6y^1\,dy^4+y^2\,dy^3)H_{\mathbf m},
 \qquad
 F_{\mathbf m,T}=-2\pi H_{\mathbf m}\eta_T.
\]
They solve the Yang--Mills equation on the flat torus.  With
\(\langle X,Y\rangle=-\tfrac12\operatorname{Tr}(XY)\),
\(D_T=\operatorname{Vol}(F_T)\), and \(S_T=6L_T-q_T\), the exact action is
\[
 \operatorname{YM}_T(\nabla^{\mathbf m,T})
 =\frac{2\pi^2k_{\mathbf m}}{g_{\rm YM}^2}
 +\frac{\pi^2s_2(\mathbf m)S_T^2}{g_{\rm YM}^2D_T}.
\]
Let \(\mathcal B_{T,\min}\le\mathcal B_{T,\max}\) be the two positive
charge-one normal-form frequencies and
\[
 \rho_T^*=\frac{\mathcal B_{T,\max}}{\mathcal B_{T,\min}},\qquad
 N_-(T)=\left\lceil\frac{\rho_T^*-1}{2}\right\rceil .
\]
For every positive root \(i<j\), write \(q_{ij}=|m_i-m_j|\).  The complete
negative spectrum of the normalized Coulomb Jacobi operator is
\[
 q_{ij}\bigl((2n+1)\mathcal B_{T,\min}
                    -\mathcal B_{T,\max}\bigr),
 \qquad 0\le n<N_-(T),
\]
and the real multiplicity contributed by this root pair is
\(12q_{ij}^2\).  Coincident numerical bands are summed.  Consequently
\[
 \operatorname{ind}_{\mathbb R}\mathcal J_{\mathbf m,T}
 =12N_-(T)\sum_{i<j}(m_i-m_j)^2
 =12rN_-(T)s_2(\mathbf m).
\]
If the distinct weights have multiplicities \(r_1,\ldots,r_\ell\), the
neutral Coulomb nullity away from a wall is
\[
 4\left(\sum_{a=1}^{\ell}r_a^2-1\right).
\]
At an odd wall \(\rho_T^*=2n+1>1\), charged zero modes add
\(12rs_2(\mathbf m)\); at \(S_T=0\), where both favourable polarizations
vanish, they add \(24rs_2(\mathbf m)\).  On the deep cusp tail,
\(N_-(T)=3\), so the index is \(36rs_2(\mathbf m)=6rk_{\mathbf m}\).
For example, \(\mathbf m=(1,0,-1)\) in \(SU(3)\) has
\[
 (q_{12},q_{13},q_{23})=(1,2,1),\qquad
 k_{\mathbf m}=12,\qquad
 \operatorname{ind}_{\mathbb R}=216,\qquad
 \operatorname{null}_{\rm neutral}=8
\]
on that tail.

These negative eigenspaces are descent directions, not bifurcation kernels;
only the wall zero modes are candidates for a Lyapunov--Schmidt analysis.
The bundle itself cannot extend over \(X\), \(S^{12}\), or \(S^{24}\):
each proposed carrier has \(H^4(-;\mathbb Z)=0\), whereas
\(c_2(\mathscr E_{\mathbf m,T})\ne0\).

\medskip\noindent
\textit{The four-dimensional \(SU(3)\) Gram quotient and its exact
boundary.}
Let
\[
 \mathscr V=\operatorname{Mat}_{3\times2}(\mathbb C),\qquad
 \Gamma(A)=A^\dagger A,
\]
with \(SU(3)\) acting by left multiplication.  Equality of Gram matrices is
exactly equality of \(SU(3)\)-orbits: after matching the two column spans by
a unitary map, the unused complex line corrects its determinant uniquely.
Thus
\[
 \mathscr V/SU(3)\cong\operatorname{Herm}_2^{\ge0}.
\]
The rank-two stratum is free and has the explicit trivialization
\[
 \mathscr V_{(2)}
 \cong SU(3)\times\operatorname{Herm}_2^{>0};
\]
the nonzero rank-one stratum has stabilizer \(SU(2)\), and the zero matrix
has stabilizer \(SU(3)\).  If \(A=(a\ b)\), then
\[
 \det\Gamma(A)
 =\lVert a\rVert^2\lVert b\rVert^2-|\langle a,b\rangle|^2
 =\lVert a\wedge b\rVert^2
 =\sum_{i<j}|a_ib_j-a_jb_i|^2.
\]
In Pauli coordinates this is
\[
 t^2-x^2-y^2-z^2,
\]
but the quotient is only the closed future cone
\(t\ge\sqrt{x^2+y^2+z^2}\), not all of Minkowski space.  Its fixed-trace
link is the Bloch ball \(B^3\), with rank-one boundary
\(\mathbb CP^1\cong S^2\).  If
\(S^{12}_{\rm col}=\mathscr V^+\) is the one-point compactification of the
representation, then
\[
 S^{12}_{\rm col}/SU(3)
 \cong\bigl(\operatorname{Herm}_2^{\ge0}\bigr)^+
 \cong B^4,
\]
and the rank-one, zero, and infinity orbits together form its boundary
\(S^3\).  This \(S^{12}_{\rm col}\) is a representation sphere; it is not
the earlier cusp-induced smash construction.

The four-variable identity
\[
 Q_4(r)=\frac12\left(\sum_{i=1}^4r_i\right)^2-\sum_{i=1}^4r_i^2
       =h^2-x^2-y^2-z^2
\]
under the orthogonal Hadamard transform is exact.  More generally,
\[
 Q_{4,\lambda}=(4\lambda-1)h^2-x^2-y^2-z^2;
\]
therefore every \(\lambda>1/4\) has Lorentz signature after rescaling time.
The coefficient \(\lambda=1/2\) is the unique choice retaining unit
coefficients in this particular orthogonal coordinate system and the
\(\pi/4\) cone angle, but it is not forced by the \(SU(3)\) orbit theorem.

The quotient does not supply a nonzero \(SU(3)\) instanton sector: its
principal bundle over \(\operatorname{Herm}_2^{>0}\cong\mathbb R^4\) is
trivial, and compactifying that open cone to \(S^4\) requires an additional
clutching class in
\(\pi_3(SU(3))\cong\mathbb Z\).  Conversely, the native torus bundle above
is topologically the pullback of an \(S^4\)-bundle along a chosen
degree-one collapse \(c_T:F_T\to S^4\), because \(c_2\) classifies
\(SU(3)\)-bundles on four-complexes.  Its native connection cannot be such
a pullback.  Indeed
\[
 \operatorname{Tr}(F_{\mathbf m,T}\wedge F_{\mathbf m,T})
 =-4\pi^2s_2(\mathbf m)\eta_T^2
\]
is nowhere zero.  If it were \(c^*\operatorname{Tr}(F_B\wedge F_B)\),
then \(dc\) would have rank four everywhere, so \(c:T^4\to S^4\) would be a
covering map, contradicting
\(\pi_1(T^4)=\mathbb Z^4\) and \(\pi_1(S^4)=0\).

\medskip\noindent
\textit{Three columns: the cubic carrier, its rank-loss sphere, and the
exact limit of the analogy.}
Now let
\[
 \mathscr V_3=\operatorname{Mat}_{3\times3}(\mathbb C).
\]
For the left \(U(3)\)-action, \(A\mapsto A^\dagger A\) is a complete orbit
invariant and gives
\[
 \mathscr V_3/U(3)
 \cong\operatorname{Herm}_3(\mathbb C)_{\geq0}.
\]
This is a nine-real-dimensional rank-three Euclidean Jordan cone.  Its
determinant is cubic, and its trace-one section is an eight-ball.  After
one-point compactification,
\[
 \mathscr V_3^+/U(3)\cong B^9,
\]
whose boundary is exactly the compactified rank-deficient locus
\[
 \{\operatorname{rank}A\le2\}\cup\{\infty\}\cong S^8.
\]
Thus the passage from two to three columns replaces the quadratic
\(3+1\) Lorentz cone and its \(S^3\) boundary by a cubic rank-three cone and
an \(S^8\) boundary at the level of the full orbit space.  The full
nine-dimensional cone is not itself a Lorentz cone; the marked angular
epimorphism nevertheless gives a split positive observation onto the
ordinary \(3+1\) cone and identifies all five retained fibre directions.

Strict \(SU(3)\) remembers one more invariant.  The complete quotient map is
\[
 A\longmapsto(H,z)=(A^\dagger A,\det A),
 \qquad |z|^2=\det H,
\]
and its image is homeomorphic to \(\mathbb R^{10}\).  Hence
\[
 S^{17}/SU(3)\cong S^9,
 \qquad
 \mathscr V_3^+/SU(3)\cong S^{10}.
\]
In this \(S^{10}\), the same compactified rank-loss set is the standard
topologically locally flat codimension-two \(S^8\), and
\[
 S^{10}\setminus S^8
 \cong\mathbb R^8\times\mathbb C^*\simeq S^1.
\]
The positive generator is the determinant phase
\(\det A/|\det A|\).  The rank-one trace-one stratum is
\(\mathbb CP^2\).  Quotienting the residual determinant-phase circle recovers
\(B^9\), with the \(S^8\) as its fixed-set boundary.  This supplies a precise
three-input carrier and alternating invariant, but not by itself a principal
bundle, connection, particle assignment, or Yang--Mills operator.

\medskip\noindent
\textit{The exact Lorentz observation and all retained directions.}
The angular map already present in the toroidal quotient,
\[
 P=\begin{pmatrix}1&-1&0\\0&1&-1\end{pmatrix},\qquad
 u=(1,1,1)^{\mathsf T},\qquad Pu=0,
\]
has the standard-Hermitian right inverse
\[
 R=P^*(PP^*)^{-1}
 =\frac13\begin{pmatrix}2&1\\-1&1\\-1&-2\end{pmatrix}.
\]
With \(\Pi=uu^*/3\) and \(Q_u=I_3-\Pi\), define
\[
 \mathcal O_P(H)=PHP^*,\qquad \jmath_P(G)=RGR^*.
\]
Both maps preserve positive semidefiniteness, and
\[
 \boxed{\mathcal O_P\jmath_P=\operatorname{id},\qquad
 \jmath_P\mathcal O_P(H)=Q_uHQ_u.}
\]
The complete linear kernel is
\[
 \ker\mathcal O_P
 =\{a\,uu^*+uv^*+vu^*:a\in\mathbb R,\ v\in u^\perp\},
\]
and its positive zero fibre is precisely
\(\{\lambda uu^*: \lambda\geq0\}\).

Set \(X=(R\;u)\).  Since
\[
 \det X=1,\qquad X^{-1}=\begin{pmatrix}P\\u^*/3\end{pmatrix},
\]
every \(H\in\operatorname{Herm}_3(\mathbb C)\) has the unique form
\[
 H=X\begin{pmatrix}G&\psi\\\psi^*&\lambda\end{pmatrix}X^*,\qquad
 G=PHP^*,\quad\psi=\frac13PHu,\quad\lambda=\frac19u^*Hu.
\]
This is the complete real decomposition
\[
 \boxed{\operatorname{Herm}_3(\mathbb C)
 \cong\operatorname{Herm}_2(\mathbb C)\oplus\mathbb C^2\oplus\mathbb R,
 \qquad 9=4+4+1,}
\]
with cubic and positive fibre
\[
 \boxed{\det H=\lambda\det G-\psi^*\operatorname{adj}(G)\psi,}
\]
\[
 H\succeq0\Longleftrightarrow
 G\succeq0,\quad\psi\in\operatorname{im}G,\quad
 \lambda\geq\psi^*G^+\psi.
\]

The embedded spin action is
\[
 T_P(g)=RgP+\Pi=X(g\oplus1)X^{-1},\qquad
 (G,\psi,\lambda)\longmapsto(gGg^*,g\psi,\lambda).
\]
For \(g\in SL_2(\mathbb C)\), the summands are a Lorentz four-vector, the
defining complex two-dimensional spin representation, and a real scalar.
The central element \(-I_2\) fixes \(G,\lambda\) and negates \(\psi\);
hence the four-vector action factors through \(SO^+(1,3)\), while the full
nine-dimensional carrier remembers the spin double cover.

The relation survives the relevant compactification in one direction:
\[
 \boxed{\jmath_P^+:B^4\cong\mathcal P_2^+
 \lhook\joinrel\longrightarrow(\partial\mathcal P_3)^+\cong S^8,}
 \qquad e_2(\jmath_P(G))=\frac13\det G,
\]
and its normal opening at \(uu^*\) is
\[
 \boxed{\det(uu^*+\varepsilon\jmath_P(G))
 =\varepsilon^2\det G.}
\]
The reverse boundary observation is nonproper.  In the strict \(SU(3)\)
quotient the determinant phase satisfies
\[
 |\omega|^2=\lambda\det G-\psi^*\operatorname{adj}(G)\psi,
\]
and \(\omega=0\) is exactly the rank-loss locus compactifying to the
distinguished \(S^8\subset S^{10}\).

The four-branch chart reaches this Lorentz sector through the holomorphic
seed
\[
 K(z)=\operatorname{diag}(z_0z_3,z_1z_2),\qquad
 \det K=z_0z_1z_2z_3.
\]
Its positive \(SL_2(\mathbb C)\)-orbit supplies every future causal
direction.  For \(Z\in\operatorname{Mat}_{3\times2}(\mathbb C)\) with the
commuting colour and spin actions, retain \(\alpha,c>0\) and put
\[
 p_\alpha(Z)=\alpha Z^*Z,\qquad
 m_{\alpha,c}^2(Z)=\frac{\alpha^2}{c^2}
 \sum_{a<b}|Z_{a1}Z_{b2}-Z_{a2}Z_{b1}|^2.
\]
The map \(p_\alpha\) separates \(SU(3)\)-orbits and is Lorentz covariant;
the mass is positive exactly at complex rank two.  Since
\[
 m_{\alpha,c}^2(rZ)=r^4m_{\alpha,c}^2(Z),
\]
 the unrestricted pointwise carrier is gapless.  Exact presymplectic reduction
now gives four distinct results.  For \(m_0>0\), the rank-two colour quotient
and its characteristic reduction are
\[
 (\mathcal J_2\times\mathcal H^+_{m_0,c})/
 ((x,p)\sim(x+sp,p))\cong_{\rm symp}T^*\mathcal H^+_{m_0,c},
\]
exactly the future-timelike spinless particle.  On nonzero rank one, the
intrinsic characteristic quotient is \(T^*\mathcal N_+^\times\), but the
written massless multiplier action is irregular because
\[
 d(\det\circ p_\alpha)_Z=0.
\]
It yields \(\dot x=\lambda p\) without \(\lambda=e\), has four transverse
rank-opening directions which are first-order null but not colour gauge, and
at \(Z=0\) admits arbitrary \(x(\tau)\).  For the separate hard constraint
with \(\omega_0\ne0\), the momentum map has rank four above its sharp
threshold and rank three at the endpoint; its action therefore forces
\(\dot x=0\) above threshold and has the ordinary
\(\dot x\in\mathbb Rp\) characteristic only at threshold.  The accompanying
multiplication/Fock spectrum is exact as a free operator calculation but is
not the Hamiltonian generating this action, so its positive endpoint is a
constraint-selected threshold rather than a dynamically generated
interacting gap.  Complete definitions, maps, stabilizers, reduced forms,
equations, and singular counterexamples are proved in the canonical module
\path{supporting_materials/workbench/research/worldline_presymplectic_reduction.tex}.

The affine normal slice is now developed as an exact field-theoretic host.
Indeed,
\[
 \mathcal N_P=\{Y\in\operatorname{Herm}_3(\mathbb C):Yu=0\}
 =\jmath_P\bigl(\operatorname{Herm}_2(\mathbb C)\bigr),
 \qquad \jmath_P^{-1}(Y)=PYP^*,
\]
and polarization of
\[
 \det(uu^*+\varepsilon\jmath_P(G))=\varepsilon^2\det G
\]
gives the constant Minkowski form of signature \((1,3)\).  Positive
semidefinite \(G\) are the future cone, rank-one positive \(G\) its null
boundary, and the lifted congruence action realizes the spin cover
\(SL_2(\mathbb C)\to SO^+(1,3)\).  The standard metrics used in the
Moore--Penrose section further select
\[
 W_P=R^*R=(PP^*)^{-1},\qquad
 K_P=\{a\in SL_2(\mathbb C):a^*W_Pa=W_P\}\cong SU(2).
\]
Its real normal representation is exactly
\(\boldsymbol1\oplus\boldsymbol3\), with trivial central action.  Promoting
this pointwise action to a local gauge symmetry gives the explicit
nonnegative Euclidean model
\[
 \int_M\left(
 \frac{|F_{\mathcal A}|^2}{2g_{\rm YM}^2}
 +\frac{Z_G}{2}|D_{\mathcal A}G|_{E_P}^2
 +\frac{m_G^2}{2}|G|_{E_P}^2
 +\frac{\lambda_4}{4}(\det G)^2\right)d\operatorname{vol}_g,
\]
where
\(E_P(G,K)=\operatorname{tr}(GW_PKW_P)\) and
\(D_{R\mathcal A P}\jmath_P(G)=\jmath_P(D_{\mathcal A}G)\).
The normal field cannot screen a fundamental centre charge; the retained
spinor \(\psi\) is centre-nontrivial, but a kinetic and interaction theory for
that additional direction has not yet been constructed.

The distinction between the candidate cusp matrix and the Gram momentum is
also now an equation:
\[
 \det G_{\eta,T}=-\frac23D_T<0,\qquad
 \operatorname{im}p_\alpha=
 \operatorname{Herm}_2(\mathbb C)_{\geq0}.
\]
Positive scaling and complex congruence preserve their different inertia
types, so any common physical model requires a new morphism rather than an
identification.  After an independent length/time dictionary is chosen, the
same Lorentzian host carries a standard Klein--Gordon/Fock model with exact
gap \(m_{\mathrm{KG}}\) for \(m_{\mathrm{KG}}>0\) and gap zero for
\(m_{\mathrm{KG}}=0\).  The massless transverse gauge Gaussian has lowest
nonzero frequency \(2\pi/L\), which closes as \(L\to\infty\), and explicit
Hamiltonian countermodels show that entanglement alone determines neither
the presence nor the size of a gap.  An auxiliary positive gap transfers to
a physical Hamiltonian only when a comparison map controls the entire
excitation form domain:
\[
 \|Ru\|\geq c_R\|u\|,\quad
 \mathfrak b[Ru]\leq C_R
 (\mathfrak a_{\rm phys}[u]-E_{\rm vac}\|u\|^2)
 \Longrightarrow
 \operatorname{gap}(A_{\rm phys})\geq\frac{\gamma c_R^2}{C_R}.
\]
An intertwiner on a single selected subspace is insufficient.  Thus a
field-theoretic kinetic operator has now been constructed on the exact
Lorentz slice, but dynamics for the complete
\((G,\psi,\lambda,\omega)\) carrier, a gauge-invariant interacting continuum
spectrum, and the comparison map required for a nonlinear mass-gap theorem
remain exact construction problems; the present theorem neither proves nor
disproves the four-dimensional Yang--Mills mass gap.

\medskip\noindent
\textit{The actual toroidal successor and its elliptic quotient.}
Independently of the Gram quotient, the rank-three Kuhn fan produces a proper
analytic degeneration
\[
 \pi_3:N_3\longrightarrow\Delta,\qquad
 N_{3,t}\cong E_\tau^3\quad(t\ne0),
\]
with cusp monodromy
\[
 M_3=\begin{pmatrix}I_3&I_3\\0&I_3\end{pmatrix}
\]
and deepest local equation \(t=z_0z_1z_2z_3\).  The irreducible central fibre
has one dense three-torus, seven double surfaces, twelve triple curves, and
six quadruple points:
\[
 (f_0,f_1,f_2,f_3)=(1,7,12,6),\qquad \chi(W_3)=6.
\]
Its normalization is the permutohedral toric threefold, with orbit vector
\[
 (1,14,36,24)\longrightarrow(1,7,12,6);
\]
the positive-codimension strata map with degrees \(2,3,4\), respectively.
Consequently, any real \(S^8\)-total-space fibration retaining \(W_3\) must
have its other singular fibres contribute Euler sum \(-4\).
For a full-rank integral shear
\(\mathcal B:\bar L_3\hookrightarrow L_3\), the calculation retains the
entire lattice index:
\[
 \chi(W_{3,\mathcal B})
 =3!\,[L_3:\mathcal B\bar L_3]
 =3!\,|\det\mathcal B|.
\]
This is the precise content of the observed \(6=3!\det I_3\).  It is a
toroidal volume/index theorem, not a proof that the Hermitian Jordan cone is
a globalization of the degeneration.

The selectively corrected Jordan determinant/Euler calculation adds an exact algebraic bridge at the
level of determinant forms.  Fixing the orientation and writing
\([a]_{\times}v=a\times v\), every \(3\times3\) Hermitian matrix is uniquely
\(S+i[a]_{\times}\), with
\[
 \boxed{\det\bigl(S+i[a]_{\times}\bigr)=\det S-a^{\mathsf T}Sa.}
\]
For a marked positive-definite symmetric monodromy form \(B\), the section
\[
 \ell_B(\lambda,a)=\lambda B+i[a]_{\times}
\]
meets the positive-semidefinite cone in exactly
\[
 \lambda\geq0,\qquad a^{\mathsf T}Ba\leq\lambda^2\det B,
\]
the future cone of a \((3+1)\)-Lorentz form.  There is also an explicit
real-linear order-cone bijection
\[
 \Theta:\operatorname{Herm}_2(\mathbb C)_{\geq0}
 \longrightarrow
 (\mathbb RI_3\oplus i\mathfrak{so}(3))
 \cap\operatorname{Herm}_3(\mathbb C)_{\geq0}
\]
such that
\[
 \det_3\Theta(H)=\frac{\operatorname{tr}H}{2}\det_2H.
\]
It sends positive-semidefinite ranks \(2,1,0\) to \(3,2,0\).  Separately,
if \(P\) is a rank-one projection in \(\operatorname{Herm}_2(\mathbb C)\),
then \(\Theta(P)\) has eigenvalues \(0,1/2,1\) and is not idempotent; hence
\(\Theta\) is not a Jordan-algebra homomorphism.  These identities exhibit
literal Lorentz sections inside the cubic cone, but
they contain no torus lattice, monodromy, period family, or filling map.  The
toric parameter tends to zero with \(B=I_3\) fixed; this is not the same as
moving to \(\det B=0\).

There is nevertheless a literal map to the reconstructed rank-two cusp.  In
the marked angular and deck lattices, set
\[
 P=\begin{pmatrix}1&-1&0\\0&1&-1\end{pmatrix},
 \qquad
 Q=\begin{pmatrix}0&-1&1\\1&-1&0\end{pmatrix},
 \qquad \Phi=P\oplus Q.
\]
Then \(\Phi:H_1(T^6)\twoheadrightarrow H_1(T^4)\) satisfies
\[
 \boxed{\Phi M_3=M_F\Phi,\qquad c_F\Phi=Qc_3.}
\]
If the candidate cusp period is \(Z_2(t)=sB_0+C(t)\), the source period
\[
 Z_3(t)=sI_3+S\,C(t)\,Q,\qquad
 S=\begin{pmatrix}1&1\\0&1\\0&0\end{pmatrix},
\]
obeys
\[
 \boxed{PZ_3(t)=Z_2(t)Q.}
\]
Consequently \(P:\mathbb C^3\to\mathbb C^2\) descends to a surjective
holomorphic map of every smooth cusp fibre, has connected elliptic kernel,
and extends over the toric boundary.  An explicit unimodular basis change
splits the source cusp integrally as
\[
 M_3\cong M_F\oplus
 \begin{pmatrix}1&1\\0&1\end{pmatrix}.
\]
The candidate's free affine order-three and order-four peripheral actions
also lift on this added elliptic factor.  The remaining problem is analytic,
not merely numerical: one still has to build a single period and overlap
family across both finite-order discs.

This successful lift is not the diagonal local-system candidate
\[
 M_1=I_3\otimes S,\qquad M_2=I_3\otimes Q,\qquad M_0=I_3\otimes T,
 \qquad M_1M_2M_0=I.
\]
Although its finite factors have orders \(4\) and \(6\), each
\(I-M_j:T^6\to T^6\) is a surjective isogeny for \(j=1,2\).  Thus every
affine lift of either finite action has a fixed point.

The six Kuhn tetrahedra also meet the Gram construction exactly at the level
of marked fans.  With
\(\nu_i=E_{ii}-I_3/3\), the \(\pi\)-th tetrahedron maps to the chamber
\[
 \mathbb R_{\ge0}\nu_{\pi_1}
 +\mathbb R_{\ge0}(-\nu_{\pi_3})
 =\{\lambda_{\pi_1}\ge\lambda_{\pi_2}\ge\lambda_{\pi_3}\}.
\]
Thus the six deepest toroidal strata and the six \(A_2\) Weyl chambers form
the same \(S_3\)-torsor after marking.  This Coxeter correspondence does not
identify the analytic germs: the toroidal germ has four branches, while the
rank-one Gram normal determinant is an irreducible quadratic.

The tempting cyclic compression of the six Kuhn labels to two also needs this
distinction.  The three even labels and three odd labels do map, through the
saturated rank-lowering homomorphism \(P:\mathbb Z^3\twoheadrightarrow
\mathbb Z^2\), to the two unimodular triangle orbits of the target fan.  They
do not arise from the asserted index-three translation superlattice: any
unimodular triangulation of a rank-three lattice fundamental domain again has
six maximal tetrahedra.

The free \(\mathbb Z^6\) belonging to this nine-dimensional-cone program is
\[
 \Lambda_{\rm tor}=H_1(T^6;\mathbb Z)
 =\mathbb Z^3_{\rm angular}\oplus\mathbb Z^3_{\rm deck}.
\]
It must not be confused either with \(\mathbb Z/6\) or with the different
rank-six bundle-defect group on the original regular locus.  Under diagonal
coordinate permutation, \(P\oplus Q\) gives
\[
 0\longrightarrow\mathbb Z(k,0)\oplus\mathbb Z(0,k)
 \longrightarrow\Lambda_{\rm tor}
 \longrightarrow A_F\oplus B_F\longrightarrow0,
 \qquad k=(1,1,1).
\]
This sequence splits as an integral lattice with monodromy, but not as an
integral \(S_3\)-module.  The quotient is two copies of the \(A_2\) weight
lattice, and the associated root lattices have index three.  The six homology
basis elements are two three-element permutation orbits; the six Kuhn/Weyl
chambers are one regular \(S_3\)-orbit.  Thus
\(\operatorname{rk}\Lambda_{\rm tor}=6\), \(|S_3|=3!=6\), and
\(\chi(W_{3,\mathcal B})=3!|\det\mathcal B|\) are not three names for one
object.  The last two share a toric cell count; the first comes from angular
plus deck homology.

The regular six-set does, however, have an exact home in the associated
nine-dimensional carrier.  Conjugation by permutation matrices gives
\[
 \operatorname{Hom}(\overline L_3,L_3)
 =\bigoplus_i\mathbb ZE_{ii}\oplus
  \bigoplus_{i\ne j}\mathbb ZE_{ij},
\]
and the six off-diagonal units form one regular \(S_3\)-orbit.  The marked
chamber \(\pi=(\pi_1\pi_2\pi_3)\) corresponds to
\(E_{\pi_1,\pi_3}\), or equivalently to the mixed angular--deck pair whose
alternating pairing is zero.  The three excluded diagonal pairs have pairing
one.  A real-linear equivariant isomorphism identifies the full matrix
representation with \(\operatorname{Herm}_3(\mathbb C)\), but the
symmetric/skew integral splitting has index \(2^3\), and the six
off-diagonal images are indefinite.  This is therefore a precise
representation/incidence bridge, not a positive, determinant, Jordan, or
toric morphism.

The only map presently reaching the positive-semidefinite Hermitian Gram
boundary is the piecewise-linear chamber map
\(\Phi_{\rm G}\circ\jmath\circ P\), which uses the angular half, is not
additive, and does not descend to a torus fibre.  Full proofs are in the
workbench sections \emph{The rank-six cusp lattice inside the
nine-dimensional cone program} and \emph{The chamber torsor inside the
nine-dimensional matrix carrier}.

\medskip\noindent
\textit{Why the higher-rung object cannot be a literal complex
fourfold \(S^8\).}
If a homotopy \(S^8\) admitted an almost-complex tangent bundle \(E\), then
\[
 c_1(E)=c_2(E)=c_3(E)=0,\qquad c_4(E)=2u,
\]
and therefore
\[
 p_2(TS^8)=c_2^2-2c_1c_3+2c_4=4u.
\]
But \(p_1=0\) and \(\operatorname{sign}(S^8)=0\), so Hirzebruch's formula
\[
 \operatorname{sign}(S^8)
 =\left\langle\frac{7p_2-p_1^2}{45},[S^8]\right\rangle
\]
forces \(p_2=0\), a contradiction.  Thus the rank-loss \(S^8\) above is a
real stratified boundary, not a smooth complex-fourfold analogue of the
candidate.  A nonsingular fibration \(T^6\to S^8\to S^2\) is also
impossible: its homotopy sequence would require a surjection
\(\mathbb Z\to\mathbb Z^6\), while Euler characteristic would require
\(2=0\).  Any real singular replacement must instead derive the actual
local filling maps; its special fibres obey
\[
 \sum_i\chi(F_i)=2.
\]
For regular-piece monodromy \(\rho\), the corrected homological Serre
filtration is
\[
 0\longrightarrow(\textstyle\bigwedge^q\Lambda)_\rho
 \longrightarrow H_q(E;\mathbb Z)
 \longrightarrow
 H_1(B^\circ;\textstyle\bigwedge^{q-1}\Lambda)
 \longrightarrow0,
\]
whereas the cohomological edge map has invariant quotient
\[
 H^q(E;\mathbb Z)\twoheadrightarrow
 (\textstyle\bigwedge^q\Lambda^\vee)^\rho.
\]
Consequently its van Kampen, Wang, and Mayer--Vietoris maps must kill the
regular fibre group and all intermediate integral homology while
respecting the coinvariant subobject and the base-one-cycle quotient;
using the reverse Serre arrows does not meet this requirement.
Every effective torus action on \(S^8\) has rank at most four, and this bound
is attained by the coordinate action
\[
 S^8/T^4\cong\operatorname{Susp}(\Delta^3)\cong B^4.
\]
This is a real two-pole singular compactification with generic \(T^4\)
orbit, not a three-special-fibre \(T^6\)-fibration over \(S^2\).

\medskip\noindent
\textit{The exact interface with the Fabel--Alp\"oge degeneration.}
The present interface does not identify the rank-loss \(S^8\) with the
candidate threefold.  It occurs on the regular locus of the reconstructed
\(S^6\).  Put
\[
 D=W\sqcup S_1\sqcup S_2,
 \qquad U=X\setminus D,
\]
where \(W\) is the non-normal central fibre and \(S_1,S_2\) are the two
multiple bielliptic fibres.  Their computed first homology and Alexander
duality give the exact map
\[
 \boxed{
 H_1(W;\mathbb Z)\oplus H_1(S_1;\mathbb Z)\oplus H_1(S_2;\mathbb Z)
 \xrightarrow{\ \sim\ }H^4(U;\mathbb Z)
 \cong\mathbb Z^2\oplus\mathbb Z^2\oplus\mathbb Z^2.}
\]
Because \(U\) has CW dimension at most five, this complete rank-six group
classifies principal \(SU(3)\)-bundles on \(U\) by \(c_2\).  The pair
sequence gives another isomorphism
\[
 \delta:H^4(U;\mathbb Z)\xrightarrow{\sim}H^5(X,U;\mathbb Z),
\]
so an unframed bundle extends across all three special fibres exactly when
its six-component class vanishes.  Global \(SU(3)\)-bundles on \(S^6\) carry
instead the separate clutching integer
\[
 k\in\pi_5(SU(3))\cong\mathbb Z,
 \qquad c_3=\pm2ku,
\]
and every such global bundle restricts trivially to \(U\).

The displayed \(\mathbb Z^6\) is a free lattice consisting of three
geometrically localized \(\mathbb Z^2\) summands.  It is not itself a cyclic
\(\mathbb Z/6\mathbb Z\).  Write its marked basis as
\[
 (a_0,b_0,a_1,b_1,a_2,b_2)
 =(\bar\gamma,\bar u,g_1,h_1,g_2,h_2).
\]
The calculations below are proved in the integrated workbench section
\emph{Fibre degree and the cyclic-six centre}, including the one-shot collar
matrix and the complete kernel, coefficient, tensor-centre, and common-centre
arguments.  They have the workbench's internal \emph{Verified} status as
lattice and matrix-group statements; that status does not extend them to a
gauge-bundle or spectral conclusion.
With the complex orientation of a smooth fibre, outward-first boundary
orientation, and the manuscript's clockwise meridians, the transpose of the
actual collar matrix gives
\[
 \operatorname{res}_F\!\left(\sum_{i=0}^2(n_i a_i+r_i b_i)\right)
 =(12n_0+4n_1-3n_2)\omega_F,
 \qquad
 [\deg_F]=(12,0,4,0,-3,0),
\]
where \(\langle\omega_F,[F]\rangle=1\).  This primitive map has saturated
kernel
\[
 \ker\deg_F
 =\mathbb Z\langle
 b_0,b_1,b_2,\ a_0-3a_1,\ 3a_1+4a_2
 \rangle.
\]
It is the pullback of the topology-killing functional, not a new numerical
coincidence.  Define
\[
 q_{\rm col}:L\longrightarrow H_1(X^\circ;\mathbb Z),\qquad
 a_0\mapsto c,\quad a_1\mapsto\bar g_1,\quad a_2\mapsto\bar g_2,\quad
 b_0,b_1,b_2\mapsto0,
\]
and let \(\Phi_1:H_1(X^\circ;\mathbb Z)\to\mathbb Z\) have row
\((12,4\ell_1,3\ell_2)\) in the basis
\((c,\bar g_1,\bar g_2)\).  Then
\[
 \deg_F=\Phi_1\circ q_{\rm col},\qquad
 \Phi_1(j_*[\widetilde a_0])
 =12\ell_0-4\ell_1-3\ell_2=p.
\]

The order-three and order-four normal-line holonomies give more than the
formal reduction of an arbitrary rank-six lattice.  If
\(\vartheta_1(a_1)=e^{-2\pi i/3}\),
\(\vartheta_2(a_2)=e^{-2\pi i/4}\), and both characters are trivial on the
other marked generators, and \(\zeta_{12}=e^{-2\pi i/12}\), then the candidate twists
\((\ell_1,\ell_2)=(1,-1)\) select
\[
 \chi_{12}(x)
 =\vartheta_1(n_1a_1+r_1b_1)
  \vartheta_2(n_2a_2+r_2b_2)^{-1}
 =\zeta_{12}^{4n_1-3n_2},\qquad
 \chi_6=\chi_{12}^2,
\]
and, for \(\rho_m(k)=e^{-2\pi i k/m}\),
\[
 \boxed{\chi_{12}=\rho_{12}\circ\deg_F,\qquad
        \chi_6=\rho_6\circ\deg_F.}
\]
Consequently
\[
 H^4(U;\mathbb Z)/\ker\chi_{12}\cong\mathbb Z/12\mathbb Z,
 \qquad
 H^4(U;\mathbb Z)/\ker\chi_6\cong\mathbb Z/6\mathbb Z.
\]
This distinguished cyclic quotient is not ordinary coefficient reduction:
\[
 H^4(U;\mathbb Z/6)\cong(\mathbb Z/6)^6.
\]
For a principal \(SU(3)\)-bundle \(P\to U\), this last identity is the exact
topological congruence
\[
 \chi_6(c_2(P))
 =\exp\!\left(-\frac{2\pi i}{6}
 \langle c_2(P)|_F,[F]\rangle\right).
\]
The full order-four normal holonomy is essential before squaring: the natural
combined character has order twelve, not six.

There is a separate exact group-theoretic occurrence of the same finite
centre.  Tensor product defines the proper closed embedding
\[
 \iota:SU(2)\times SU(3)\hookrightarrow SU(6),\qquad
 \iota(A,B)=A\otimes B,
\]
and its restriction
\(\mu_2\times\mu_3\xrightarrow{\sim}\mu_6=Z(SU(6))\) sends
\((\lambda,\mu)\) to \(\lambda\mu\).  Since \(U\) has homotopy dimension at
most five, \(c_2\) classifies \(SU(n)\)-bundles for \(n\ge3\).  If
\(e=c_2(E)\) and \(f=c_2(F_3)\) for determinant-trivial bundles of ranks two
and three, then the splitting principle gives
\[
 c_2(E\otimes F_3)=3e+2f.
\]
Hence every \(SU(6)\)-bundle class \(h\in H^4(U;\mathbb Z)\) admits at least
one tensor-subgroup reduction, for example \((e,f)=(h,-h)\); all solutions at
the \(c_2\)-level are
\[
 (e,f)=(h+2t,-h-3t),\qquad t\in H^4(U;\mathbb Z),
\]
so neither a nonzero \(h\) nor a reduction is canonically selected.
Moreover
\[
 H^5(U;\mathbb Z)
 \cong\ker\{\mathbb Z^3\xrightarrow{(1,1,1)}\mathbb Z\}
\cong Q(A_2),
\]
after equipping the distinguished three-component support basis with the
combinatorial orthonormal form.  The resulting simple-root Gram matrix is
\(\bigl(\begin{smallmatrix}2&-1\\-1&2\end{smallmatrix}\bigr)\); it is not a
cup-product or intersection form.  This support-redistribution lattice is not
the defect lattice and does not make its rank-five fibre-degree kernel an
\(A_5\) root lattice.

The isomorphism
\(H^4(U)/\ker\chi_6\cong Z(SU(6))\) is an isomorphism of finite groups, not a
system of \(SU(6)\) transition functions or a connection.  In particular,
the injective tensor map does not factor through
\((SU(2)\times SU(3))/\mu_6\): its central \(\mu_6\) maps to nontrivial
scalar matrices.  The distinct common-centre construction is
\[
 \Gamma_6=\{((\lambda\mu)^{-1},\lambda I_2,\mu I_3):
              \lambda\in\mu_2,\ \mu\in\mu_3\},
 \qquad
 \bigl(U(1)\times SU(2)\times SU(3)\bigr)/\Gamma_6
 \hookrightarrow U(6),\qquad [z,A,B]\longmapsto z(A\otimes B),
\]
where \(\Gamma_6\cong\mu_6\).  Its determinant is \(z^6\), and its image
meets \(SU(6)\) in exactly \(\iota(SU(2)\times SU(3))\); the image in
\(U(6)\) is not generally contained in \(SU(6)\).  What remains
unproved is a degree-changing morphism identifying this degree-four number
with the degree-two lifting obstruction of a common-centre quotient gauge
bundle and with its electric--magnetic line pairing.  No mass-gap statement
follows before that bundle, a connection, operator, domain, and spectral
comparison have been constructed.

\medskip\noindent
\textit{The audited Erd\H{o}s--Straus selector and its Koide boundary
model.}
Let \(P\equiv1\pmod4\) be prime and let
\(\mathscr S(P)\) be the set of sorted positive solutions of
\[
 \frac4P=\frac1x+\frac1y+\frac1z.
\]
Every solution has a unique least denominator
\[
 x=a_j=\frac{P+R_j}{4},\qquad R_j=4j+3,\qquad
 0\le j<\frac{P-1}{4}.
\]
For \(a=a_j\), every divisor \(u\mid a^2\) belongs to one of two tagged
selector channels,
\[
 E_j:\ R_j\mid4u+1,\qquad
 M_j:\ R_j\mid4u+P.
\]
For \(\varepsilon=0\) on \(E_j\) and \(\varepsilon=1\) on \(M_j\), put
\[
 T=\frac{4u+P^\varepsilon}{R_j},\qquad
 K=\frac{P^{1-\varepsilon}T+1}{4}.
\]
Then
\[
 \frac4P
 =\frac1a+
  \frac1{P^\varepsilon aK/u}+
  \frac1{PK}.
\]
Middle-channel complementation
\(\sigma(u)=a^2/u\) exchanges the last two denominators and has no fixed
point.  The corrected selector theorem is the bijection
\[
 \mathcal E(P)\ \sqcup\
 \bigl(\mathcal M(P)/\langle\sigma\rangle\bigr)
 \xrightarrow{\ \sim\ }\mathscr S(P).
\]
The missing point in the supplied proof was that \(a_j\) remains the unique
least denominator: if another denominator \(b<a_j\), then
\(1/a_j+1/b>2/a_j\ge4/P\), contradicting positivity of the third summand.

Writing \(v=a^2/u\) gives the positive rank-one matrix
\[
 \Gamma(a,u)=\begin{pmatrix}u&a\\a&v\end{pmatrix}.
\]
It has the unique primitive factorization
\[
 u=hr^2,\qquad a=hrs,\qquad v=hs^2,\qquad \gcd(r,s)=1,
\]
and the two selector congruences become
\[
 E_j:\ R_j\mid Pr+s,\qquad
 M_j:\ R_j\mid r+s.
\]
Consequently the unresolved prime case of the Erd\H{o}s--Straus conjecture
is equivalently the assertion that for every prime \(P\equiv1\pmod{12}\)
there exist \(h,r,s>0\) with
\[
 \gcd(r,s)=1,\qquad
 \frac P4<hrs<\frac P2,\qquad
 R=4hrs-P,
\qquad
 R\mid Pr+s\ \ \text{or}\ \ R\mid r+s.
\]
This is an exact reformulation; it does not prove the required nonemptiness.

The accompanying terminal matrix is also exact.  With
\[
 \ell_0=\frac{T+1}{4},\qquad \ell_1=u-jT,\qquad \ell_2=1,
\]
\[
 \lambda=12\ell_0-4\ell_1-3\ell_2=P^\varepsilon,\qquad
 \mathsf R=
 \begin{pmatrix}3&-\ell_1\\4&1-4\ell_0\end{pmatrix},\qquad
 \det\mathsf R=-P^\varepsilon,
\]
and \(\operatorname{SNF}(\mathsf R)=\operatorname{diag}(1,P^\varepsilon)\).
The pair \((\Gamma,\mathsf R)\) is a faithful arithmetic certificate,
whereas \(\mathsf R\) alone has explicit collisions.  For exterior states
\(\lambda=+1\) and \(\det\mathsf R=-1\); the actual reconstructed
\(S^6\) candidate uses
\[
 (\ell_0,\ell_1,\ell_2)=(0,1,-1),\qquad
 \lambda=-1,\qquad \det\mathsf R=+1.
\]
No \(SL_2(\mathbb Z)\) row operation identifies these two oriented
presentations, since it preserves the signed determinant.

Finally,
\[
 G(t,m,b)=\begin{pmatrix}2t&2m\\2m&b/3\end{pmatrix},\qquad
 (t,m,b)=\left(\frac u2,\frac a2,3v\right)
\]
gives \(G=\Gamma\) and
\[
 \det G=
 \left(\frac{u+v}{2}\right)^2-a^2-
 \left(\frac{u-v}{2}\right)^2=0.
\]
The inverse orthogonal Koide transform therefore produces a real Koide-null
vector.  It lies in the positive three-coordinate chamber only when, for
\(\theta_K=u/a=r/s\),
\[
 0<\theta_K<\frac{\sqrt3-\sqrt2}{1+\sqrt2}
 \quad\text{or}\quad
 \frac{\sqrt3+\sqrt2}{1+\sqrt2}<\theta_K<1+\sqrt2.
\]
Most selector states are not automatically positive mass points.

This coefficient correspondence lies on the discriminant of the ambient
period-matrix space:
\[
 \det P_{\rm selector}=tb-6m^2=0,
\]
whereas a regular torus fibre requires
\(\det P_T=-D_T\ne0\), with \(\det P_T\sim-T^2\) on the radial cusp.
Thus the displayed arithmetic map cannot be a map to a regular fibre.
Together with the signed terminal mismatch, this proves the exact current
boundary: the arithmetic incidence and Koide/Lorentz identities are real,
but they do not define a geometric morphism to the candidate threefold.
The complete proofs and correction ledger are in
\path{supporting_materials/workbench/research/} and the
Erd\H{o}s--Straus bridge audit supplied with the workbench.

The shared residue-lattice continuation supplies an additional exact
count map, without altering these specific geometric boundaries.
For $p\equiv1\pmod{24}$ prime and $R>0$, $R\equiv3\pmod4$, retain
\[
 a=(p+R)/4,\qquad g=\gcd(R,pa),\qquad R_0=R/g,\qquad N_0=pa/g.
\]
The complete divisor-to-ordered-pair bijection is
\[
 d\mid N_0^2,\quad d>0,\quad d\equiv-N_0\pmod{R_0}
 \quad\longleftrightarrow\quad
 (b,c)=\left(\frac{N_0+d}{R_0},\frac{N_0+N_0^2/d}{R_0}\right),
\]
where the target consists of all positive integer pairs with
$4/p=1/a+1/b+1/c$, and the inverse is $d=R_0b-N_0$.
The raw count $A_R(p)$ and unordered count satisfy
$U_R(p)=(A_R(p)+\mathbf1_{R_0\mid2})/2$.
For a finite set $P$ of primes reducing into $S_{840}=\langle289\rangle$,
put $w_{R,P}(r)=\sum_{p\in P,\ p\equiv r\ (840)}A_R(p)$,
$W_{R,P}=\sum_r w_{R,P}(r)e_r$, $M_{R,P}=\sum_r w_{R,P}(r)$,
and $u=\sum_{r\in S_{840}}e_r$. The exact integral map is
\[
 Q_{R,P}=6W_{R,P}-M_{R,P}u\in A(S_{840})\subset N.
\]
The shared source's ``Inversion and an integral augmentation-valued
count map'' theorem gives the full Fourier inverse with the fixed
generator $289$. Both $R$ and the sample $P$ remain part of this map;
it is neither a universal positivity proof nor a geometric map to $X$.

\medskip\noindent
\textit{Exact Koide-shell parametrization, hyperbolic realization, and the
remaining analytic problem.}
The arithmetic count isolated by the reconstruction is
\[
 R(X)=\#\left\{(a,q,n)\in\mathbb Z^3:
 \begin{array}{c}
 a^2-3q^2-18n^2=-2,\quad a>0,\ q,n\geq0,\\
 a+4n>q,\quad 4a+17n\leq X
 \end{array}\right\}.
\]
It now has a multiplicity-one integral parametrization.  Put
\[
 \mathcal O_K=\mathbb Z[\sqrt3],\qquad
 \epsilon=2+\sqrt3,\qquad m_n=18n^2-2\quad(n\geq1).
\]
The ideals of norm $m_n$ are in bijection with
\[
 \mathcal P(m_n)=\{(C,D,E):C,D>0,\ C^2D=m_n,\ 0\leq E<D,
                         \ E^2\equiv3\pmod D\}
\]
by
\[
 (C,D,E)\longmapsto C\bigl(D\mathbb Z+(E+\sqrt3)\mathbb Z\bigr).
\]
For every such triple there is one reduced generator
\[
 C(A+B\sqrt3),\qquad
 A^2-3B^2=D,\quad A\equiv EB\pmod D,\quad A>2B\geq0,
\]
and all nonnegative-$q$ points in its orbit are obtained, exactly once, from
\[
 a_{n,C,D,E,k}+q_{n,C,D,E,k}\sqrt3
   =C(A+B\sqrt3)\epsilon^k,\qquad k\geq0.
\]
The inverse is equally integral: form the HNF data of
$(a+q\sqrt3)$, then apply
\[
 (a,q)\longmapsto(2a-3q,\,2q-a)
\]
while $a\leq2q$.  The number of steps is $k$, and the terminal pair is
the unique reduced generator.  The omitted face is not hidden in this
description: for $n=0$ the solutions are precisely
\[
 a_k+q_k\sqrt3=(1+\sqrt3)\epsilon^k,\qquad k\geq1.
\]

Writing
\[
 \rho(m)=\sum_{d\mid m}\left(\frac{12}{d}\right),
\]
the number of unit orbits at fixed $n\geq1$ is
\[
 \#\mathcal P(m_n)=\rho(m_n)
 =\rho(3n-1)\rho(3n+1).
\]
The original sharp height is retained rather than replaced by a radial norm.
For the reduced generator $\alpha_{n,\mathbf p}$, set
\[
 \tau_{n,\mathbf p}=
 \frac{\log(\alpha_{n,\mathbf p}/\sqrt{m_n})}{\log\epsilon}\in[0,1),
 \qquad Y_n(X)=\frac{X-17n}{4},
\]
\[
 L_n(X)=\frac1{\log\epsilon}
 \log\!\left(
 \frac{Y_n(X)+\sqrt{Y_n(X)^2-m_n}}{\sqrt{m_n}}
 \right),
 \quad
 \mathcal N(X)=\{n\geq1:Y_n(X)\geq\sqrt{m_n}\}.
\]
Then the endpoint-exact formula is
\[
 R(X)=R_0(X)+
 \sum_{n\in\mathcal N(X)}\ \sum_{\mathbf p\in\mathcal P(m_n)}
 \bigl(\lfloor L_n(X)-\tau_{n,\mathbf p}\rfloor+1\bigr),
\]
where $R_0(X)$ is the preceding $n=0$ Pell count.  Thus the strict
chamber wall and every equality $4a+17n=X$ remain visible.  Fourier
expansion of the phases gives unitary Hecke coefficients
\[
 S_\ell(m)=\sum_{N\mathfrak a=m}\xi_\ell(\mathfrak a),\qquad
 S_\ell(18n^2-2)=(-1)^\ell S_\ell(3n-1)S_\ell(3n+1),
\]
so the unproved remainder is an explicit sharply truncated shifted
correlation, rather than an unspecified equidistribution assumption.

There is also an exact quaternionic realization.  In
\[
 B=(18,-6)_{\mathbb Q}
 =\mathbb Q\langle I,J\rangle/(I^2-18,\ J^2+6,\ IJ+JI),
\]
the lattice
\[
 \mathcal O=\mathbb Z1\oplus
 \mathbb Z\frac{I+J}{2}\oplus
 \mathbb Z\frac{I-J}{2}\oplus\mathbb Z\frac{IJ}{2}
\]
is an order of reduced discriminant $108$, of index $18$ in a maximal
overorder.  The algebra has discriminant $6$, ramifies exactly at $2$
and $3$, and splits over $\mathbb R$.  The map
\[
 (a,q,n)\longmapsto
 \gamma=2+\frac{qI+aJ+nIJ}{2}
\]
is a bijection from the parity-restricted ternary shell to the locus
\[
 \{\gamma\in\mathcal O:\operatorname{trd}(\gamma)=4,
                           \ \operatorname{nrd}(\gamma)=1\}.
\]
Every image is primitive hyperbolic, has eigenvalues $2\pm\sqrt3$, and
has translation length $2\log(2+\sqrt3)$.  Its centralizer order is
\[
 \mathcal O\cap\mathbb Q(\gamma)=
 \mathbb Z[\gamma-2]\cong\mathbb Z[\sqrt3].
\]
The full shell is a finite union of conjugacy classes under
$\Gamma=\mathcal O^1$, and $R(X)$ is exactly the sum over their left
cosets by the centralizers subject to the same four chamber inequalities and
the same height $4a+17n\leq X$.  No factor two is introduced on passing to
$PSL_2(\mathbb R)$.

This realization explains both what is now proved and what is still open.
The standard radial Huber parameter is
\[
 X_i(\gamma)^2=1+\frac{(7a+15\sqrt2\,n)^2}{48},
\]
not the Koide height $4a+17n$.  The two heights are comparable but their
sharp regions have a main-order symmetric difference, so a radial remainder
does not transfer to the original count.  Benoist--Oh's effective theorem
does apply directly to the literal five-wall sector and proves, for some
$\eta>0$,
\[
 R(X)=C_0X+O(X^{1-\eta}),\qquad
 C_0=\frac9{2\pi^2}\log(2+\sqrt3)
                 \log\!\left(\frac3{2\sqrt2}\right).
\]
The strongest audited nearby pointwise sharp radial theorem has exponent
$2/3$; the smooth radial spectral estimate has exponent $39/64$, but
neither theorem has the target sector and height.  Therefore
\[
 R(X)-C_0X=O_\varepsilon(X^{1/2+\varepsilon})
\]
remains \textbf{OPEN}, not a polite substitute for a theorem: its exact
missing object is the displayed sharp shifted Hecke correlation, including
all endpoint and exceptional-spectrum terms.

The exact census independently verifies every point through $X=10^8$:
\[
 R(10^8)=3{,}536{,}002.
\]
For integer $10^4\leq X\leq10^8$, the largest observed value of
$|R(X)-C_0X|/\sqrt X$ is
$0.1714001928088145\ldots$, attained at $X=867{,}240$; this is finite
evidence for the open estimate, not its proof.  These arithmetic and
hyperbolic results do not by themselves identify the counted triples with
physical lepton masses or a quantum-field-theoretic spectrum.  They instead
give the exact morphisms and analytic target that any such identification
must respect.

\end{document}
